\chapter{Where It All Meets}

\begin{center}\itshape
Beauty is the first test: there is no permanent place in the world for ugly mathematics.\\
\upshape --- G.~H.~Hardy, \emph{A Mathematician's Apology} (1940), \S10
\end{center}

In March 1948, at the Institute for Advanced Study in Princeton, Atle Selberg found a formula about the primes that used nothing deeper than the
Möbius function and Mertens's estimates. ``I found the fundamental formula \dots{} in March this year'', he wrote to Hermann Weyl that September. In April he
squeezed out of it a curious consequence: if $a$ and $A$ are the lower and upper limits of $\frac{\theta(x)}x$, then $a+A=2$. The prime number theorem says
$a=A=1$, and Selberg saw that it would follow from either half. In early July Paul Turán, who had asked for Selberg's notes on another proof, mentioned the formula in
a seminar; Paul Erd\H{o}s was in the audience. Within days Erd\H{o}s had deduced from it that for every $\varepsilon>0$ the interval $]x,(1+\varepsilon)x]$ contains
at least a fixed proportion of $\frac x{\log x}$ primes. ``Sunday afternoon (July~18) I used his result'', Selberg recalled fifty years later: that day he had
an elementary proof of the prime number theorem, the first that used no complex analysis. What followed was less elementary. The two men
could not agree on whether to publish together, and in 1949 the proofs appeared separately, Selberg's in the \emph{Annals of Mathematics} and
Erd\H{o}s's in the \emph{Proceedings of the National Academy of Sciences}. Selberg received the Fields Medal in 1950. Dorian Goldfeld has told the story from the
letters of both sides.

Selberg's formula is a meeting point in the purest sense: the Möbius function of Tome~1, the hyperbola method of Tome~5, Chapter~2, and Mertens's theorems of Tome~5,
Chapter~1, combine into a statement none of them suggests alone. Problem~\ref{pb:selberg} proves it. This last chapter of the book is made of such meetings. Every problem
is a Journey: a long climb in many steps, each within reach, ending at a theorem with a name. The four sections follow four threads of Part~I. The first is about
squares and divisor sums, where Jacobi's theta function, the $p$-adic numbers, Landau's Euler products and Ramanujan's Eisenstein series meet. The second listens to the zeros of $\zeta$,
which know where the primes are, and to the primes, which know where the zeros are. The third puts characters to work, from primitive roots to Euler's famous
polynomial $n^2+n+41$. The last mixes angles, chance and cotangents.

\begin{think}
Euler noticed that $n^2+n+41$ is prime for $n=0,1,\dots,39$. Guess what fraction of the values for $n<10^6$ are prime, and compare with the fraction of primes among
integers of the same size. A second guess: is $\pi(x)$ larger or smaller than $\operatorname{li}(x)$ for $x\le10^8$, and how often does Riemann's function
$\Pi(x)=\pi(x)+\frac12\pi(x^{1/2})+\frac13\pi(x^{1/3})+\cdots$ cross $\operatorname{li}(x)$ in that range?
\end{think}

Throughout, $p$ and $q$ denote primes, $\Lambda$ is von Mangoldt's function, $\theta(x)=\sum_{p\le x}\log p$, $\psi(x)=\sum_{n\le x}\Lambda(n)$, $\sigma(n)$ and
$\sigma_k(n)$ are sums of divisors and of their $k$-th powers, $\omega(n)$ is the number of prime factors of $n$, $\chi$ is a Dirichlet character, $e(t)=e^{2\pi it}$, and
$\operatorname{li}(x)=\int_0^x\frac{\mathrm dt}{\log t}$ (principal value), so that $\operatorname{li}(x)=\operatorname{Li}(x)+1.0451\ldots$ with $\operatorname{Li}$ as in
Tome~5. Earlier problems of this book are used freely; in particular Problem~\ref{pb:ramanujan5} proved Jacobi's triple product
\[
\prod_{n\ge1}(1-x^{2n})(1+zx^{2n-1})(1+z^{-1}x^{2n-1})=\sum_{m\in\Z}x^{m^2}z^m\qquad(|x|<1,\ z\ne0)
\]
and Jacobi's identity $\prod_{n\ge1}(1-x^n)^3=\sum_{k\ge0}(-1)^k(2k+1)x^{k(k+1)/2}$.

\section{Squares and divisor sums}

Jacobi proved in 1829 that the number of ways to write $n$ as a sum of four squares is eight times the sum of the divisors of $n$ not divisible by $4$. The four journeys of
this section come at such counts from four directions: from $q$-series, from congruences modulo prime powers, from Euler products, and from modular forms. Each
direction proves something the others cannot see.

\begin{exercise}[Journey]\label{pb:jacobi}
\emph{Jacobi's four-square theorem.} For $|x|<1$ let $\Theta(x)=\sum_{m\in\Z}x^{m^2}$, so that $\Theta(x)^4=\sum_{n\ge0}r_4(n)x^n$, where $r_4(n)$ is the number of
$(a,b,c,d)\in\Z^4$ with $a^2+b^2+c^2+d^2=n$. Write $E(x)=\prod_{n\ge1}(1-x^n)$, and $D=x\frac{\mathrm d}{\mathrm dx}$.
\begin{enumerate}[label=(\alph*)]
\item Show that $\Theta(-x)=\prod_{n\ge1}\frac{1-x^n}{1+x^n}$.
\item Show that $E(x)^6=\frac12\sum_{r,s\in\Z}\bigl((2r+1)^2-4s^2\bigr)x^{r^2+r+s^2}$.
\item Let $B(x)=\sum_{r\in\Z}x^{r^2+r}$. Show that $2E^6=\Theta\,(B+4DB)-4B\,D\Theta$, and that
\[
\Theta(x)=\prod_{n\ge1}(1-x^{2n})(1+x^{2n-1})^2,\qq B(x)=2\prod_{n\ge1}(1-x^{2n})(1+x^{2n})^2.
\]
\item Deduce that
\[
\Theta(-x)^4=1-8\sum_{n\ge1}\Bigl(\frac{(2n-1)x^{2n-1}}{1+x^{2n-1}}-\frac{2nx^{2n}}{1+x^{2n}}\Bigr),
\]
and then that $r_4(n)=8\sum_{d\mid n,\,4\nmid d}d$ for every $n\ge1$.
\item Deduce Lagrange's theorem that every positive integer is a sum of four squares, and compute $r_4(p)$ for an odd prime $p$ and $r_4(2^k)$.
\item By counting lattice points in a ball of $\R^4$, show that $\sum_{n\le x}r_4(n)=\frac{\pi^2}2x^2+O(x^{3/2})$. Show also, from (d), that
$\sum_{n\le x}r_4(n)=3\zeta(2)x^2+O(x\log x)$, and conclude that $\zeta(2)=\frac{\pi^2}6$. (The notebook finds $\sum_{n\le10^7}r_4(n)=493\,480\,252\,693\,888$, which is
$1.00000007\cdot\frac{\pi^2}2\,10^{14}$.)
\end{enumerate}
\end{exercise}

\begin{exercise}[Journey]\label{pb:densities}
\emph{Four squares, one prime at a time.} For a prime $p$, $k\ge0$ and $n\ge1$, let $N_{p^k}(n)$ be the number of $x\in(\Z/p^k\Z)^4$ with
$x_1^2+x_2^2+x_3^2+x_4^2\equiv n\pmod{p^k}$ (so $N_1(n)=1$). Call such an $x$ \emph{primitive} if some $x_i$ is not divisible by $p$, and let $N^*_{p^k}(n)$ count the
primitive ones.
\begin{enumerate}[label=(\alph*)]
\item Let $p$ be odd. Show that $N^*_{p^{k+1}}(n)=p^3N^*_{p^k}(n)$ for $k\ge1$, and that for $k\ge2$, $N_{p^k}(n)-N^*_{p^k}(n)$ is $p^4N_{p^{k-2}}(\frac n{p^2})$ if
$p^2\mid n$, and $0$ otherwise.
\item Deduce that the \emph{local density} $\delta_p(n)=\lim_{k\to\infty}p^{-3k}N_{p^k}(n)$ exists, and, using Problem~\ref{pb:sumsquares}, that
$\delta_p(n)=\bigl(1-\frac1{p^2}\bigr)\bigl(1+\frac1p+\dots+\frac1{p^a}\bigr)$ when $p^a$ is the exact power of $p$ dividing $n$.
\item Let $p=2$. Show that $N^*_{2^{k+1}}(n)=8N^*_{2^k}(n)$ for $k\ge3$. (Every lift of a solution modulo $2^k$ has the same sum of squares modulo $2^{k+1}$; pair the
solutions.) Deduce from a count modulo $8$ that $\delta_2(n)=1$ for odd $n$, and $\delta_2(n)=3\cdot2^{-a}$ when $2^a$ exactly divides $n$ and $a\ge1$.
\item Show that for every $n\ge1$,
\[
r_4(n)=\pi^2n\prod_p\delta_p(n),
\]
the product converging absolutely. Interpret the factor $\pi^2n$ as a density at infinity: the derivative of the volume of the ball $x_1^2+\dots+x_4^2\le n$.
\end{enumerate}
\end{exercise}

\begin{exercise}[Journey]\label{pb:twosquares}
\emph{How many integers are sums of two squares?} Let $b(n)=1$ if $n\ge1$ is a sum of two squares and $b(n)=0$ otherwise, $B(x)=\sum_{n\le x}b(n)$, and
$F(s)=\sum_{n\ge1}b(n)n^{-s}$ for $s>1$. Let $\chi_4$ be the non-principal character modulo $4$.
\begin{enumerate}[label=(\alph*)]
\item Show that for $s>1$,
\[
F(s)^2=\frac{\zeta(s)L(s,\chi_4)}{1-2^{-s}}\prod_{p\equiv3\,(4)}\frac1{1-p^{-2s}},
\]
and deduce that $F(\sigma)\sim C(\sigma-1)^{-1/2}$ as $\sigma\to1^+$, with $C^2=\frac\pi2\prod_{p\equiv3\,(4)}(1-p^{-2})^{-1}$.
\item Let $S(y)=\sum_{n\le y}\frac{b(n)}n$ and $\delta=\frac1{\log y}$. Show that $S(y)\le e\,F(1+\delta)$ and $S(y)\ge F(1+\delta)-e^{-1/2}F(1+\frac\delta2)$, and
deduce that $S(y)$ lies between two positive multiples of $\sqrt{\log y}$ for large $y$.
\item Show that $\sum_{p^k\le y}k\log p=O(y)$ and $\sum_{n\le x}b(n)\log n\le c\,x\,S(x)$ for an absolute constant $c$, and deduce that $B(x)=O\bigl(\frac x{\sqrt{\log x}}\bigr)$.
\item Show that $\sum_{n\le x}b(n)\log n\ge\sum_{m\le x/\log^3x}b(m)\,\theta(\frac xm;4,1)$, where $\theta(y;4,1)=\sum_{p\le y,\,p\equiv1\,(4)}\log p$, and
deduce that $\frac{B(x)\sqrt{\log x}}x$ is bounded below by a positive constant.
\item Show that \emph{if} $B(x)\sim\frac{Kx}{\sqrt{\log x}}$ for some constant $K$, then $K=\frac C{\sqrt\pi}=0.7642\ldots$ The notebook finds $B(10^8)=18\,457\,847$, for which
$\frac{B(x)\sqrt{\log x}}x=0.7922$.
\end{enumerate}
\end{exercise}

\begin{exercise}[Journey]\label{pb:ramanujanE2}
\emph{Ramanujan's differential equations.} With $q=e^{2\pi i\tau}$ for $\tau$ in the upper half-plane, let $E_2=1-24\sum_{n\ge1}\sigma(n)q^n$, $E_4=1+240\sum_{n\ge1}\sigma_3(n)q^n$,
$E_6=1-504\sum_{n\ge1}\sigma_5(n)q^n$, $\eta(\tau)=q^{1/24}\prod_{n\ge1}(1-q^n)$, and $D=\frac1{2\pi i}\frac{\mathrm d}{\mathrm d\tau}=q\frac{\mathrm d}{\mathrm dq}$.
\begin{enumerate}[label=(\alph*)]
\item Show that $\frac{D\eta}\eta=\frac{E_2}{24}$.
\item Using Dedekind's transformation $\eta(-\frac1\tau)=\sqrt{-i\tau}\,\eta(\tau)$, show that $E_2(-\frac1\tau)=\tau^2E_2(\tau)+\frac{6\tau}{\pi i}$.
\item Show that $f=DE_2-\frac{E_2^2}{12}$ satisfies $f(-\frac1\tau)=\tau^4f(\tau)$, and deduce \emph{Ramanujan's equations}
\[
DE_2=\frac{E_2^2-E_4}{12},\qq DE_4=\frac{E_2E_4-E_6}3,\qq DE_6=\frac{E_2E_6-E_4^2}2.
\]
\item Deduce that for every $n\ge1$,
\[
\sum_{m=1}^{n-1}\sigma(m)\sigma(n-m)=\frac{5\sigma_3(n)+(1-6n)\sigma(n)}{12}.
\]
\item Deduce similarly that $\sum_{m=1}^{n-1}\sigma(m)\sigma_3(n-m)=\frac{21\sigma_5(n)+(10-30n)\sigma_3(n)-\sigma(n)}{240}$.
\item For a prime $p$, show that $\sum_{m=1}^{p-1}\sigma(m)\sigma(p-m)=\frac{(p-1)(5p^2-p-6)}{12}$, and check $p=7$ by hand.
\end{enumerate}
\end{exercise}

\section{What the zeros know}

Riemann's explicit formula (Tome~5, Theorem~7.5) writes the primes as a sum over the zeros of $\zeta$. Read from left to right, it says that the zeros determine the primes;
read from right to left, that the primes determine the zeros. The first two journeys make both readings audible. The third asks what Cramér's random model of Tome~5,
Chapter~10, gets wrong about the sign of $\pi(x)-\operatorname{li}(x)$, and the fourth proves Selberg's formula, the starting point of the elementary proof of the
prime number theorem, without any zero at all.

\begin{exercise}[Journey]\label{pb:landauzeros}
\emph{The zeros hear the primes.} Fix a real number $x>1$, and let $\Lambda(x)=\Lambda(n)$ if $x$ is an integer $n$, and $\Lambda(x)=0$ otherwise. The zeros
$\rho=\beta+i\gamma$ of $\zeta$ in the critical strip are counted with multiplicity.
\begin{enumerate}[label=(\alph*)]
\item Show that for $T\ge2$,
\[
\frac1{2\pi}\int_1^Tx^{2+it}\,\frac{\zeta'}\zeta(2+it)\dd t=-\frac T{2\pi}\Lambda(x)+O_x(1).
\]
\item Using $\xi(s)=\xi(1-s)$ (Tome~5, Proposition~6.2), show that
$\frac{\zeta'}\zeta(s)=-\frac{\zeta'}\zeta(1-s)+\log\pi-\frac12\frac{\Gamma'}\Gamma\bigl(\frac s2\bigr)-\frac12\frac{\Gamma'}\Gamma\bigl(\frac{1-s}2\bigr)$, and deduce that
$\int_1^Tx^{-1+it}\frac{\zeta'}\zeta(-1+it)\dd t=O_x(\log T)$.
\item Deduce \emph{Landau's formula}: for $T\ge3$,
\[
\sum_{0<\gamma\le T}x^\rho=-\frac T{2\pi}\Lambda(x)+O_x(\log^2T).
\]
\item Assume the Riemann hypothesis. Show that $\sum_{0<\gamma\le T}\cos(\gamma\log x)=-\frac T{2\pi}\frac{\Lambda(x)}{\sqrt x}+O_x(\log^2T)$: the frequencies $\gamma$
resonate exactly at the prime powers. (Figure~\ref{fig:landau} shows it with the first $2000$ zeros.)
\item Still under the Riemann hypothesis, let $0<\gamma_1\le\gamma_2\le\cdots$ be the positive ordinates. Show that for every real $\alpha\ne0$ the sequence
$(\alpha\gamma_n)$ is equidistributed modulo $1$.
\end{enumerate}
\end{exercise}

\begin{figure}[ht]
\centering
\fig{ch14-landau}
\caption{The zeros hear the primes: $-\frac{2\pi\sqrt x}T\sum_{0<\gamma\le T}\cos(\gamma\log x)$ for $1.5\le x\le30$, with the first $2000$ zeros
($T=\gamma_{2000}$). By Problem~\ref{pb:landauzeros}(d) it tends to $\Lambda(x)$: spikes of height $\log p$ (red dots) at the prime powers $p^k$, and nothing in
between.}\label{fig:landau}
\end{figure}

\begin{exercise}[Journey]\label{pb:hearzeros}
\emph{The primes hear the zeros.} Fix $s>0$ and $T\ge0$, let $h(t)=e^{-(t-T)^2/2s^2}+e^{-(t+T)^2/2s^2}$, and
\[
g(u)=\frac1{2\pi}\int_{-\infty}^\infty h(t)e^{-iut}\dd t=\frac{2s}{\sqrt{2\pi}}\cos(Tu)\,e^{-s^2u^2/2}.
\]
\begin{enumerate}[label=(\alph*)]
\item Let $H(z)=h(-i(z-\frac12))$. Show that $H$ is entire, that $H(1-z)=H(z)$, that $H(\frac12+it)=h(t)$, and that $H(\sigma+it)$ tends to $0$ faster than any
power of $|t|$, uniformly for $\sigma$ in a bounded interval.
\item Integrate $H(z)\frac{\zeta'}\zeta(z)$ around rectangles with vertices $-1\pm iT'$, $2\pm iT'$ to prove the \emph{Guinand--Weil explicit formula}
\[
\sum_\rho H(\rho)=H(0)+H(1)-2\sum_{n\ge1}\frac{\Lambda(n)}{\sqrt n}g(\log n)-g(0)\log\pi+\frac1{2\pi}\int_{-\infty}^\infty h(t)\operatorname{Re}\frac{\Gamma'}\Gamma\Bigl(\frac14+\frac{it}2\Bigr)\dd t,
\]
the sum running over all non-trivial zeros, and converging absolutely.
\item Assume the Riemann hypothesis, and let $Z(T)=\sum_{\gamma>0}e^{-(\gamma-T)^2/2s^2}$, a sum of bumps of width $s$, one at each zero. Show that for $T\ge1$,
\[
Z(T)=\frac s{\sqrt{2\pi}}\log\frac T{2\pi}-\frac{2s}{\sqrt{2\pi}}\sum_{n\ge2}\frac{\Lambda(n)}{\sqrt n}\cos(T\log n)\,e^{-s^2\log^2n/2}+O_s\Bigl(\frac1T\Bigr).
\]
\item The right side uses only the primes. With $s=0.6$ each term with $n>e^{10}$ is below $10^{-9}$, and all of them together are below $10^{-6}$. Figure~\ref{fig:hearzeros} plots it for $5\le T\le50$. Read off
the first zeros. Why does a smaller $s$ separate the zeros better, and what does it cost?
\end{enumerate}
\end{exercise}

\begin{figure}[ht]
\centering
\fig{ch14-hearzeros}
\caption{The primes hear the zeros: the right side of Problem~\ref{pb:hearzeros}(c) with $s=0.6$, computed from the prime powers $n\le e^{10}$ alone (blue). Its peaks
sit at the ordinates of the zeros of $\zeta$ (red ticks: $14.13$, $21.02$, $25.01$, $30.42$, $32.94$, $37.59$, $40.92$, $43.33$, $48.01$, $49.77$).}\label{fig:hearzeros}
\end{figure}

\begin{exercise}[Journey]\label{pb:cramersign}
\emph{The coin and the sign of $\pi(x)-\operatorname{li}(x)$.} Let $\pi_{\mathcal C}(x)=\sum_{n\le x}X_n$ be the counting function of Cramér's random primes
(Tome~5, Definition~10.1), and let $\Pi(x)=\sum_{k\ge1}\frac1k\pi(x^{1/k})$ be Riemann's prime-power counting function.
\begin{enumerate}[label=(\alph*)]
\item Show that $\mathbb E\,\pi_{\mathcal C}(x)=\operatorname{li}(x)+O(1)$ and $\operatorname{Var}\pi_{\mathcal C}(x)\sim\frac x{\log x}$, and that
$P\bigl(\pi_{\mathcal C}(x)>\operatorname{li}(x)\bigr)\to\frac12$.
\item Show that, almost surely, $\pi_{\mathcal C}(x)-\operatorname{li}(x)$ changes sign infinitely often. (Use the independent blocks $]n_{j-1},n_j]$ with $n_j=2^{2^j}$, and
both Borel--Cantelli lemmas.)
\item Show that $\pi(x)=\Pi(x)-\frac12\pi(\sqrt x)+O(x^{1/3})$, and that
\[
\Pi(x)-\operatorname{li}(x)=\frac{\psi(x)-x}{\log x}+\int_2^x\frac{\psi(t)-t}{t\log^2t}\dd t+\frac2{\log2}-\operatorname{li}(2).
\]
\item The notebook finds, for $10^3\le x\le10^8$: $\pi(x)<\operatorname{li}(x)$ throughout, with $\frac{\pi(x)-\operatorname{li}(x)}{\sqrt x/\log x}$ between $-2.55$ and
$-0.53$; while $\Pi(x)-\operatorname{li}(x)$ changes sign $66\,112$ times below $10^8$, with $\frac{\Pi(x)-\operatorname{li}(x)}{\sqrt x/\log x}$ between $-0.75$ and $0.71$. In the
model, $\pi_{\mathcal C}(x)-\operatorname{li}(x)$ is typically of size $\sqrt{x/\log x}$, which is $2330$ at $x=10^8$. Explain the sign of $\pi(x)-\operatorname{li}(x)$
with (c). Then explain, with (c) and the explicit formula $\psi(x)=x-\sum_\rho\frac{x^\rho}\rho+\cdots$ of Tome~5, Theorem~7.5, why the true fluctuations are of
size $\frac{\sqrt x}{\log x}$ rather than $\sqrt{x/\log x}$ if the Riemann hypothesis holds.
\end{enumerate}
\end{exercise}

\begin{exercise}[Journey]\label{pb:selberg}
\emph{Selberg's formula.} Let $\Lambda_2=\mu*\log^2$, that is, $\Lambda_2(n)=\sum_{d\mid n}\mu(d)\log^2\frac nd$.
\begin{enumerate}[label=(\alph*)]
\item Show that $\Lambda_2(n)=\Lambda(n)\log n+\sum_{dm=n}\Lambda(d)\Lambda(m)$. In particular $\Lambda_2\ge0$, and $\Lambda_2(n)=0$ when $n$ has three distinct prime
factors.
\item Show that $\sum_{d\le x}\frac{\mu(d)}d\log\frac xd=O(1)$ and $\sum_{d\le x}\frac{\mu(d)}d\log^2\frac xd=2\log x+O(1)$.
\item Deduce that $\sum_{n\le x}\Lambda_2(n)=2x\log x+O(x)$.
\item Deduce \emph{Selberg's formula}
\[
\theta(x)\log x+\sum_{pq\le x}\log p\,\log q=2x\log x+O(x),
\]
the sum running over ordered pairs of primes.
\item Let $a=\liminf\frac{\theta(x)}x$ and $A=\limsup\frac{\theta(x)}x$. Show that $a+A=2$.
\item The notebook computes $\frac1x\bigl(\theta(x)\log x+\sum_{pq\le x}\log p\log q-2x\log x\bigr)$: it is $-3.668$, $-3.667$, $-3.668$, $-3.666$ at $x=10^4$, $10^5$,
$10^6$, $10^7$, although the left side of the formula is only $0.886$ times $2x\log x$ at $10^7$. Explain why so slow an approach to $2x\log x$ is consistent with an
error term $O(x)$.
\end{enumerate}
\end{exercise}

\section{Characters at work}

Dirichlet invented characters to find primes in arithmetic progressions (Tome~5, Chapter~4). They turned out to be good at almost everything: detecting primitive
roots, counting residues, sorting primes by their angles. In this section they meet the Fourier matrix, binary quadratic forms, Brun--Titchmarsh and the
Pólya--Vinogradov inequality of Tome~5, Theorem~11.11, which says that $\bigl|\sum_{M<n\le M+N}\chi(n)\bigr|\le\sqrt q\,(1+\log q)$ for every primitive character
$\chi$ modulo $q$.

\begin{exercise}[Journey]\label{pb:primroot}
\emph{The least primitive root.} Let $p$ be an odd prime, let $g(p)$ be the least primitive root modulo $p$, and for $1\le N<p$ let $P(N)$ be the number of primitive
roots modulo $p$ in $\{1,\dots,N\}$.
\begin{enumerate}[label=(\alph*)]
\item Show that for every $n$ prime to $p$,
\[
\mathbf 1_{\text{$n$ is a primitive root}}=\sum_{e\mid p-1}\frac{\mu(e)}e\sum_{\chi^e=\chi_0}\chi(n),
\]
the inner sum running over the $e$ characters modulo $p$ whose $e$-th power is the principal character $\chi_0$.
\item Show that $\bigl|P(N)-\frac{\varphi(p-1)}{p-1}N\bigr|<2^{\omega(p-1)}\sqrt p\,(1+\log p)$, and that the same bound holds for the number of primitive roots in any
interval of $N$ consecutive integers, $N<p$.
\item Deduce that $g(p)<2^{\omega(p-1)}\frac{p-1}{\varphi(p-1)}\sqrt p\,(1+\log p)+1$, and that $g(p)=O(p^{1/2+\varepsilon})$ for every $\varepsilon>0$.
\item Deduce that for every $\varepsilon>0$, every interval of length $p^{1/2+\varepsilon}$ contains about the expected proportion $\frac{\varphi(p-1)}{p-1}$ of
primitive roots, once $p$ is large.
\item The record values of $g(p)$ for $p<10^6$ are
\[
g(3)=2,\ g(7)=3,\ g(23)=5,\ g(41)=6,\ g(71)=7,\ g(191)=19,\ g(409)=21,\ g(2161)=23,
\]
\[
g(5881)=31,\ g(36721)=37,\ g(55441)=38,\ g(71761)=44,\ g(110881)=69,\ g(760321)=73,
\]
and for $1000<p<10^6$, $g(p)$ never exceeds $0.0055$ times the bound of (c). Which part of the bound is wasteful?
\end{enumerate}
\end{exercise}

\begin{exercise}[Journey]\label{pb:schur}
\emph{Schur's determinant and the sign of the Gauss sum.} Let $p$ be an odd prime, $\zeta=e(\frac1p)$, $\chi$ the Legendre symbol modulo $p$, and
$G=\sum_{n=0}^{p-1}\zeta^{n^2}$. Let $F$ be the $p\times p$ matrix $(\zeta^{jk})_{0\le j,k<p}$.
\begin{enumerate}[label=(\alph*)]
\item Show that $G=\sum_{n=1}^{p-1}\chi(n)\zeta^n$, that $G^2=\chi(-1)p$, and that $\operatorname{tr}F=G$.
\item Show that $F^2=pR$, where $R$ is the permutation matrix of $j\mapsto-j$ modulo $p$. Deduce that the eigenvalues of $F$ lie in $\{\sqrt p,-\sqrt p,i\sqrt p,-i\sqrt p\}$,
and that their multiplicities $a,b,c,d$ satisfy $a+b+c+d=p$, $a+b-c-d=1$ and $G=\sqrt p\,\bigl(a-b+i(c-d)\bigr)$.
\item Writing $F$ as a Vandermonde matrix, show that $\det F=i^{p(p-1)/2}\,\Delta$ with $\Delta>0$.
\item Comparing with $\det F=\sqrt p^{\,p}\,(-1)^b\,i^{\,c-d}$, show that $G=\sqrt p$ if $p\equiv1\pmod4$ and $G=i\sqrt p$ if $p\equiv3\pmod4$.
\item Let $p\equiv3\pmod4$, $S=\sum_{0<a<p/2}\chi(a)$, $U=\sum_{0<a<p/2}a\chi(a)$ and $T=\sum_{a=1}^{p-1}a\chi(a)$. Show that $T=2U-pS$ and $T=\chi(2)(4U-pS)$, and,
using Dirichlet's formula for $L(1,\chi)$ (Tome~5, Theorem~3.10) together with (d), that
\[
S=\bigl(2-\chi(2)\bigr)\frac{\sqrt p}\pi\,L(1,\chi).
\]
Deduce that for $p\equiv3\pmod4$ there are more quadratic residues than non-residues in $]0,\frac p2[$, and that $L(1,\chi)\ge\frac\pi{3\sqrt p}$.
\item The notebook finds $(S,L(1,\chi))=(1,\,1.18741)$, $(3,\,1.96520)$, $(3,\,0.246069)$, $(39,\,1.27940)$, $(77,\,2.41818)$ for $p=7$, $23$, $163$, $1019$, $10007$.
Check the formula of (e) on these values, and explain why $p=163$ gives exactly $L(1,\chi)=\frac\pi{\sqrt{163}}$.
\end{enumerate}
\end{exercise}

\begin{exercise}[Journey]\label{pb:euler41}
\emph{Euler's polynomial $n^2+n+41$.} A positive definite binary quadratic form $[a,b,c]=ax^2+bxy+cy^2$ has discriminant $D=b^2-4ac<0$ and $a>0$. Two forms are
\emph{equivalent} if one is obtained from the other by a substitution $(x,y)\mapsto(\alpha x+\beta y,\gamma x+\delta y)$ with $\alpha\delta-\beta\gamma=1$ and integer
coefficients. A form is \emph{reduced} if $|b|\le a\le c$, with $b\ge0$ when $|b|=a$ or $a=c$.
\begin{enumerate}[label=(\alph*)]
\item Show that every positive definite form is equivalent to a reduced form of the same discriminant, that a reduced form has $a\le\sqrt{|D|/3}$, and that equivalent
forms represent the same integers with coprime $(x,y)$.
\item Let $q$ be a prime, $f(n)=n^2+n+q$ and $D=1-4q$. Prove \emph{Rabinowitsch's theorem}: $f(n)$ is prime for $n=0,1,\dots,q-2$ if and only if $[1,1,q]$ is the only
reduced form of discriminant $D$. (If $p<q$ divides $f(n)$, look at $[p,2n+1,\frac{f(n)}p]$.)
\item Find the reduced forms of discriminant $-163$, and deduce Euler's observation. The notebook finds that the primes $q\le5000$ with the property of (b) are $2$, $3$,
$5$, $11$, $17$, $41$.
\item Let $\chi(p)=\bigl(\frac{-163}p\bigr)$. Show that for an odd prime $p$ the number of solutions of $f(n)\equiv0\pmod p$, for $q=41$, is $1+\chi(p)$, and that $f(n)$ is
never even. The \emph{Hardy--Littlewood conjecture} predicts that $\#\{n<N:f(n)\text{ prime}\}\sim C\sum_{n<N}\frac1{\log f(n)}$, with
\[
C=2\prod_{p>2}\Bigl(1-\frac{\chi(p)}{p-1}\Bigr).
\]
Explain where $C$ comes from, and show that $C=\frac43\cdot\frac{\sqrt{163}}\pi\prod_{p>2}\frac{1-\chi(p)/(p-1)}{1-\chi(p)/p}$, a product that converges much faster.
The notebook finds $C=6.6395$, and $261\,081$ primes among $f(n)$, $n<10^6$, against the prediction $261\,020$.
\item For $n^2+1$ the same heuristic gives the constant $C'=\prod_{p>2}\bigl(1-\frac{\chi_4(p)}{p-1}\bigr)=1.3728$, and the notebook counts $54\,110$ primes among
$n^2+1$, $1\le n<10^6$, against the prediction $53\,972$. Explain in words why $n^2+n+41$ is almost five times richer in primes than $n^2+1$.
\end{enumerate}
\end{exercise}

\begin{exercise}[Journey]\label{pb:erdosnp}
\emph{The least quadratic non-residue, on average.} For an odd prime $p$ let $n(p)$ be the least quadratic non-residue modulo $p$; it is a prime. Let
$p_1=2,p_2=3,p_3=5,\dots$ be the primes in order.
\begin{enumerate}[label=(\alph*)]
\item Show that for fixed $k$, the proportion of primes $p\le x$ with $n(p)=p_k$ tends to $2^{-k}$. (By quadratic reciprocity, the condition is a condition on $p$ modulo
$M_k=8p_2\cdots p_k$.)
\item Let $y\ge2$ with $8\prod_{3\le q\le y}q\le\sqrt x$. Using the Brun--Titchmarsh inequality in progressions (Tome~5, Exercise~8.4), show that
$\#\{p\le x:n(p)>y\}\le c\,\frac{\pi(x)}{2^{\pi(y)}}$ for an absolute constant $c$.
\item Deduce that
\[
\frac1{\pi(x)}\sum_{p\le x}\min\Bigl(n(p),\frac{\log x}3\Bigr)\longrightarrow E=\sum_{k\ge1}\frac{p_k}{2^k}=3.67464\ldots
\]
\item The notebook finds that the untruncated mean $\frac1{\pi(x)}\sum_{3\le p\le x}n(p)$ is $3.4943$, $3.6004$, $3.6461$, $3.6631$, $3.6706$ for $x=10^4,\dots,10^8$, that
the proportions of $n(p)=2,3,5,7,11,13,17$ below $10^8$ are $0.500061$, $0.250012$, $0.125011$, $0.062525$, $0.031265$, $0.015576$, $0.007857$, and that the largest value is
$n(48\,473\,881)=67$. Below $10^8$, $719\,697$ primes have $n(p)>\frac{\log x}3=6.14$, and the truncated mean of (c) is only $3.1422$. Explain this last number, and
explain what more would be needed to prove that the untruncated mean tends to $E$.
\end{enumerate}
\end{exercise}

\section{Angles, chance and cotangents}

The last four journeys leave the beaten track. Hecke's characters sort the Gaussian primes by their angles, the coin of Part~I is tossed for the prime factors of
$p-1$, a sieve meets the powers of $2$, and the sawtooth function of Tome~4 is taken apart with cotangents and residues.

\begin{exercise}[Journey]\label{pb:hecke}
\emph{The angles of the Gaussian primes.} Every prime $p\equiv1\pmod4$ is $p=a^2+b^2$ with unique integers $a>b>0$ (Tome~1, Theorems~7.3 and~10.4). Let
$\varphi_p=\arctan\frac ba\in\,]0,\frac\pi4[$. \emph{Hecke's theorem} says that the angles $\varphi_p$ are equidistributed in $]0,\frac\pi4[$.
\begin{enumerate}[label=(\alph*)]
\item For $k\in\Z$ and $\alpha\in\Z[i]\setminus\{0\}$ let $\lambda_k(\alpha)=\bigl(\frac\alpha{|\alpha|}\bigr)^{4k}$. Show that $\lambda_k(\alpha)$ depends only on the ideal
$(\alpha)$ and that $\lambda_k(\alpha\beta)=\lambda_k(\alpha)\lambda_k(\beta)$. Let $S_k(x)=\sum_{0<|\alpha|^2\le x}\lambda_k(\alpha)$. Show that $S_k(x)=O\bigl((1+|k|)\sqrt x\bigr)$
for $k\ne0$, by comparing each term with the average of $\lambda_k$ over a unit square.
\item Let $L_k(s)=\frac14\sum_{\alpha\ne0}\lambda_k(\alpha)|\alpha|^{-2s}$. Show that $L_k$ is holomorphic on $\operatorname{Re}s>\frac12$ for $k\ne0$, that
$L_0(s)=\zeta(s)L(s,\chi_4)$, and that for $\operatorname{Re}s>1$
\[
L_k(s)=\frac1{1-(-1)^k2^{-s}}\prod_{p\equiv1\,(4)}\frac1{(1-e^{4ik\varphi_p}p^{-s})(1-e^{-4ik\varphi_p}p^{-s})}\prod_{p\equiv3\,(4)}\frac1{1-p^{-2s}}.
\]
\item Using $3+4\cos\theta+\cos2\theta\ge0$, show that $L_k(1+it)\ne0$ for all real $t$ when $k\ne0$.
\item Let $\vartheta_k(x)=\sum_{p\le x,\ p\equiv1\,(4)}2\cos(4k\varphi_p)\log p$. Using Newman's theorem (Tome~5, Theorem~5.3) and Tome~5, Lemma~5.4, applied to the
non-decreasing functions $\vartheta_0+\vartheta_k$ and $\vartheta_0-\vartheta_k$, show that $\vartheta_k(x)=o(x)$ for $k\ne0$.
\item Deduce Hecke's theorem. (Only cosines appear: reflect the angles $4\varphi_p\in\,]0,\pi[$ to $]-\pi,0[$ and use Weyl's criterion, Tome~4, Theorem~6.13.)
\item Among the $332\,180$ primes $p\equiv1\pmod4$ below $10^7$, the notebook finds a proportion $0.499795$ with $\varphi_p\le\frac\pi8$, and a proportion $0.12658$
with $b<\frac a{10}$. What does Hecke's theorem predict for the second number? Figure~\ref{fig:angles} shows the histogram.
\end{enumerate}
\end{exercise}

\begin{figure}[ht]
\centering
\fig{ch14-angles}
\caption{Angles $\varphi_p=\arctan\frac ba$ of the $332\,180$ primes $p=a^2+b^2\equiv1\pmod4$ below $10^7$, in $20$ bins, each bar divided by the count expected for the
uniform distribution. All bars lie between $0.991$ and $1.007$.}\label{fig:angles}
\end{figure}

\begin{exercise}[Journey]\label{pb:shifted}
\emph{How many prime factors has $p-1$?} Let $\omega(n)$ be the number of distinct prime factors of $n$, and for $y\ge2$ let $\omega_y(n)$ count only those $\le y$.
\begin{enumerate}[label=(\alph*)]
\item Fix $y$. Show that the law of $\omega_y(p-1)$, for $p$ chosen uniformly among the primes $\le x$, tends as $x\to\infty$ to the law of $\sum_{q\le y}B_q$, where the
$B_q$ are independent and $P(B_q=1)=\frac1{q-1}=1-P(B_q=0)$.
\item Show that $\sum_{q\le y}\frac1{q-1}=\log\log y+M+\sum_q\frac1{q(q-1)}+o(1)$, where $M=0.26149\ldots$ is the constant in Mertens's theorem (Tome~5, Theorem~1.6),
so that the constant is $1.03465\ldots$, and that $\frac{\sum_{q\le y}B_q-\log\log y}{\sqrt{\log\log y}}$ tends in law to $\mathcal N(0,1)$ as $y\to\infty$.
\item Using the Brun--Titchmarsh inequality (Tome~5, Exercise~8.4), show that $\sum_{p\le x}\omega(p-1)\le(2+o(1))\,\pi(x)\log\log x$.
\item Assuming the Bombieri--Vinogradov theorem (stated after Tome~5, Proposition~8.11), show that $\frac1{\pi(x)}\sum_{p\le x}\omega(p-1)=\log\log x+O(1)$.
\item The notebook finds, over the primes $p\le x$ for $x=10^4,\dots,10^8$: mean of $\omega(p-1)$ equal to $3.027$, $3.298$, $3.518$, $3.696$, $3.847$; mean minus
$\log\log x$ equal to $0.807$, $0.855$, $0.892$, $0.916$, $0.933$; variance $0.540$, $0.656$, $0.781$, $0.906$, $1.021$. Below $10^8$ the proportions of
$\omega(p-1)=2,3,4,5,6,7$ are $0.0835$, $0.2922$, $0.3709$, $0.2041$, $0.0459$, $0.0033$. For comparison, $\omega(n)$ for $n\le10^8$ has mean $3.150$ and variance $1.214$,
while the model of (a) with all $q\le10^8$ has mean $3.948$ and variance $2.573$. Why might one expect the mean to approach $\log\log x+1.035$, and why is the variance
so much smaller than in the model?
\end{enumerate}
\end{exercise}

\begin{exercise}[Journey]\label{pb:romanoff}
\emph{Romanoff's theorem.} Let $r(n)$ be the number of pairs $(p,k)$, $p$ prime and $k\ge1$, with $n=p+2^k$, and let $\mathcal R=\{n:r(n)>0\}$.
\begin{enumerate}[label=(\alph*)]
\item Show that $\sum_{n\le x}r(n)\ge c\,x$ for some $c>0$ and all large $x$, and that $\#(\mathcal R\cap[1,x])\ge\frac{(\sum_{n\le x}r(n))^2}{\sum_{n\le x}r(n)^2}$.
\item Using the upper bound $\#\{p\le x:p+d\text{ prime}\}\le c\,\bigl(\frac d{\varphi(d)}\bigr)^2\frac x{\log^2x}$ for even $d\ge2$, which follows from Selberg's sieve as in
Tome~5, Exercise~8.5, show that $\sum_{n\le x}r(n)^2=O(x)$ provided $\sum_{m\le L}f(m)=O(L)$, where $f(m)=\bigl(\frac{2^m-1}{\varphi(2^m-1)}\bigr)^2$.
\item Show that $\bigl(\frac n{\varphi(n)}\bigr)^2=\sum_{e\mid n}\mu(e)^2g(e)$, where $g$ is multiplicative with $g(q)=\frac{2q-1}{(q-1)^2}\le\frac5q$ for odd primes $q$, and
deduce that $\sum_{m\le L}f(m)\le L\sum_{e\text{ odd}}\frac{\mu(e)^2g(e)}{\operatorname{ord}_e(2)}$.
\item Show that every odd squarefree $e$ with $\operatorname{ord}_e(2)\le J$ divides $P_J=\prod_{j\le J}(2^j-1)<2^{J^2}$, deduce that
$\sum_{\operatorname{ord}_e(2)\le J}\mu(e)^2g(e)=O\bigl((\log J)^5\bigr)$, and conclude \emph{Romanoff's theorem}: $\mathcal R$ has positive lower density.
\item Let $n\equiv7\,629\,217\pmod{11\,184\,810}$. Check that $11\,184\,810=2\cdot3\cdot5\cdot7\cdot13\cdot17\cdot241$ and that for every $k\ge1$ the number $n-2^k$ is
divisible by one of $3,5,7,13,17,241$. (Look at $k$ modulo $2,4,3,12,8,24$.) Deduce that infinitely many odd numbers are not in $\mathcal R$.
\item The notebook finds that the proportion of integers $n\le x$ in $\mathcal R$ is $0.475$, $0.4662$, $0.4605$, $0.4578$ for $x=10^4,\dots,10^7$, and among odd $n$
it is $0.947$, $0.932$, $0.921$, $0.916$, while $\frac1L\sum_{m\le L}f(m)=2.155$, $2.372$, $2.487$, $2.535$ for $L=10$, $20$, $40$, $80$. Comment.
\end{enumerate}
\end{exercise}

\begin{exercise}[Journey]\label{pb:dedekind}
\emph{Dedekind sums.} Let $((t))=t-\lfloor t\rfloor-\frac12$ for $t\notin\Z$ and $((t))=0$ for $t\in\Z$. For coprime integers $h$ and $k\ge1$, the \emph{Dedekind sum} is
$s(h,k)=\sum_{r=0}^{k-1}\bigl(\bigl(\frac rk\bigr)\bigr)\bigl(\bigl(\frac{hr}k\bigr)\bigr)$.
\begin{enumerate}[label=(\alph*)]
\item Show that $\bigl(\bigl(\frac rk\bigr)\bigr)=-\frac1{2k}\sum_{j=1}^{k-1}\cot\frac{\pi j}k\,\sin\frac{2\pi jr}k$ for every integer $r$, and deduce that
\[
s(h,k)=\frac1{4k}\sum_{j=1}^{k-1}\cot\frac{\pi j}k\cot\frac{\pi hj}k.
\]
\item Let $h,k\ge1$ be coprime. By integrating $\cot(\pi z)\cot(\pi hz)\cot(\pi kz)$ around the rectangle with vertices $\varepsilon\pm iY$, $1+\varepsilon\pm iY$, and letting
$Y\to\infty$, prove \emph{Dedekind's reciprocity law}
\[
s(h,k)+s(k,h)=-\frac14+\frac1{12}\Bigl(\frac hk+\frac kh+\frac1{hk}\Bigr).
\]
\item Show that $s(1,k)=\frac{(k-1)(k-2)}{12k}$, that $6k\,s(h,k)$ is an integer, and that $|s(h,k)|\le s(1,k)$.
\item Let $\frac hk=[0;a_1,\dots,a_n]$ with $0<h<k$. Show that $|s(h,k)|\le\frac1{12}(a_1+\dots+a_n)+n$, and that $s(h,k)$ can be computed in $O(\log k)$ arithmetic
operations on fractions.
\item For $k=101$, $1009$, $10\,007$, $100\,003$ the notebook finds that the average of $|s(h,k)|$ over $1\le h<k$ is $0.951$, $2.037$, $3.691$, $5.802$, while $\log^2k$ is
$21.3$, $47.8$, $84.8$, $132.5$, and that the maximum is $s(1,k)$. Vardi proved in 1993 that $\frac{s(h,k)}{\log k}$ has a Cauchy limiting distribution as $h$ and $k$
vary. Explain why a law with no mean is consistent with averages growing like $\log^2k$ rather than $\log k$.
\end{enumerate}
\end{exercise}

\section{Hints}

\begin{hint}{pb:jacobi}
(a) Put $z=-1$ in the triple product and use $\prod(1+x^n)\prod(1-x^{2n-1})=1$. (b) Multiply Jacobi's identity for $E^3$ by itself and pair $k$ with $l$ through $r=\frac{k+l}2$
or $\frac{k+l-1}2$. (c) Write $E^6$ as a product of a theta series in $x^2$ and a series in $x$, using the triple product with $x\mapsto x^2$ twice; or verify the identity
by differentiating the logarithms. (d) Logarithmic derivatives of the products in (a) and (c). (f) Each point of $\Z^4$ in the ball owns the unit cube in front of it.
\end{hint}

\begin{hint}{pb:densities}
(a) Hensel: a primitive solution has some $x_i$ prime to $p$, and $x_i\mapsto x_i+p^kt$ changes the sum by $2x_ip^kt$. A non-primitive solution is $px'$. (b) Use $N_p(n)$
from Problem~\ref{pb:sumsquares}, and induct on $a$. (c) Odd squares are $1$ modulo $8$. (d) Compare the product with $8\sum_{d\mid n,4\nmid d}d$ prime by prime.
\end{hint}

\begin{hint}{pb:twosquares}
(a) Compare Euler factors, using Tome~1, Theorem~7.3. (b) For $n\le y$, $\frac1n\le e\,n^{-1-\delta}$; for the lower bound, the tail $n>y$ is at most
$y^{-\delta/2}F(1+\frac\delta2)$. (c) Write $\log n=\sum_{p^k\parallel n}k\log p$; removing $p^k$ from $n$ keeps $b=1$. For $B(x)$, compare $\sum b(n)\log n$ with
$\log\frac x{\log x}\,\bigl(B(x)-B(\frac x{\log x})\bigr)$. (d) $m\cdot p$ with $p\equiv1\pmod4$ is a sum of two squares when $m$ is. (e) Abel's theorem for Dirichlet series: if $B(x)\sim\frac{Kx}{\sqrt{\log x}}$,
compute $F(\sigma)$ by partial summation.
\end{hint}

\begin{hint}{pb:ramanujanE2}
(a) Logarithmic derivative of the product. (b) Differentiate $\log\eta(-\frac1\tau)=\log\eta(\tau)+\frac12\log(-i\tau)$. (c) $f$ is holomorphic at $i\infty$ with $f=-\frac1{12}+\dots$;
use the dimensions $\dim M_4=\dim M_6=\dim M_8=1$ (Tome~5, Theorem~12.9) for $f$, for $DE_4-\frac{E_2E_4}3$ and for $DE_6-\frac{E_2E_6}2$. (d) Compare coefficients of $q^n$ in
$12DE_2=E_2^2-E_4$.
\end{hint}

\begin{hint}{pb:landauzeros}
(a) Expand $\frac{\zeta'}\zeta(2+it)=-\sum\Lambda(n)n^{-2-it}$ and integrate term by term. (b) Take the logarithmic derivative of the functional equation, and integrate by
parts against $x^{it}$. (c) Integrate $x^s\frac{\zeta'}\zeta(s)$ around the rectangle $-1\le\sigma\le2$, $1\le t\le T'$ with a good height $T'$ (Tome~5, Lemma~7.3). (e) Weyl's criterion
with $\alpha\gamma_n$; $\sum_{\gamma\le T}e^{2\pi ik\alpha\gamma}$ is the sum of (d) with $x=e^{2\pi k\alpha}$, and $N(T)\sim\frac T{2\pi}\log T$.
\end{hint}

\begin{hint}{pb:hearzeros}
(a) $h$ extends to an entire even function of $t$. (b) The rectangle encloses the zeros with $|\gamma|<T'$ and the pole at $1$; on $\operatorname{Re}s=-1$ use the
logarithmic derivative of the functional equation, and shift the $\Gamma$ term to the line $\frac12$, picking up the pole of $\frac{\Gamma'}\Gamma(\frac s2)$ at $s=0$. (c) Under the
Riemann hypothesis $H(\rho)=h(\gamma)$. Use $\operatorname{Re}\frac{\Gamma'}\Gamma(\frac14+\frac{it}2)=\log\frac{|t|}2+O(t^{-2})$.
\end{hint}

\begin{hint}{pb:cramersign}
(a) $\sum_{3\le n\le x}\frac1{\log n}-\operatorname{li}(x)$ is bounded; the Lindeberg condition is trivial for bounded variables. (b) With $V_j=\operatorname{Var}\pi_{\mathcal C}(n_j)$, the
block increment exceeds $\operatorname{li}$'s increment by $2\sqrt{V_j}$ with probability bounded below, while $|\pi_{\mathcal C}(n_{j-1})-\operatorname{li}(n_{j-1})|\le\sqrt{V_j}$
eventually. (c) $\psi(x)=\sum_k\theta(x^{1/k})$; for the identity, integrate $\frac{d\psi}{\log t}$ by parts. (d) $\psi(x)-x$ is about $\sqrt x$ times a bounded
oscillating sum, and $\pi(\sqrt x)\approx\frac{2\sqrt x}{\log x}$.
\end{hint}

\begin{hint}{pb:selberg}
(a) Differentiate $\Lambda=\mu*\log$ with the derivation $f\mapsto f\log$, which satisfies $(f*g)\log=(f\log)*g+f*(g\log)$. (b) Write $\log\frac xd$ and $\log^2\frac xd$ through
$\sum_{m\le x/d}$ of simple functions, and use $\sum_{d\mid n}\mu(d)=[n=1]$ and $\sum_{m\le y}\frac{\tau(m)}m=\frac12\log^2y+2\gamma\log y+c+O(y^{-1/2}\log y)$. (c) $\sum_{n\le x}\Lambda_2(n)=\sum_d\mu(d)\sum_{m\le x/d}\log^2m$.
(e) Apply Selberg's formula at $x$ and at $\frac xp$, and average.
\end{hint}

\begin{hint}{pb:primroot}
(a) $\sum_{\chi^e=\chi_0}\chi(n)$ is $e$ or $0$ according as $n$ is an $e$-th power. (b) Isolate the principal character and use Pólya--Vinogradov for each other one.
(c) If $N=g(p)-1$, then $P(N)=0$. Use $2^{\omega(m)}=O(m^\varepsilon)$ and $\frac m{\varphi(m)}=O(\log\log m)$.
\end{hint}

\begin{hint}{pb:schur}
(b) $(F^2)_{jk}=\sum_l\zeta^{l(j+k)}$. $\operatorname{tr}F^2=p$. (c) $\zeta^k-\zeta^j=2i\,e\bigl(\frac{j+k}{2p}\bigr)\sin\frac{\pi(k-j)}p$. (d) $(a-b)^2+(c-d)^2=1$ by (a); go through the
four cases. (e) Pair $a$ with $p-a$, and split $T$ by the parity of $a$.
\end{hint}

\begin{hint}{pb:euler41}
(a) Repeatedly replace $[a,b,c]$ by $[a,b-2ma,\dots]$ (translation) or by $[c,-b,a]$ (rotation); $a$ decreases. (b) For $\Leftarrow$, a reduced form $[a,2n+1,c]$ with $a>1$ gives
$ac=f(n)$. For $\Rightarrow$, the principal form does not represent $p<q$. (d) Count roots of $(2n+1)^2\equiv-163\pmod p$. Dirichlet's formula for $L(1,\chi)$ gives the
product $\prod(1-\chi(p)/p)^{-1}$.
\end{hint}

\begin{hint}{pb:erdosnp}
(a) The vector $\bigl((\frac{p_i}p)\bigr)_{i\le k}$ is determined by $p$ modulo $M_k$, and each sign pattern takes the same number of classes. (b) Sum the Brun--Titchmarsh
bound over the classes with all signs $+1$. (c) Split at a fixed $k_0$, and bound the rest with (b).
\end{hint}

\begin{hint}{pb:hecke}
(a) A unit $i^j$ has $(i^j)^{4k}=1$. Across a unit square $\lambda_k$ varies by $O(\frac{|k|}{|\alpha|})$, and $\int_0^{2\pi}e^{4ik\theta}\dd\theta=0$. (c) Look at
$|L_0(\sigma)^3L_k(\sigma+it)^4L_{2k}(\sigma+2it)|\ge1$. (d) The Laplace transform of $\vartheta_k$ is a combination of $\frac{L_k'}{L_k}$ and $\frac{L_{-k}'}{L_{-k}}$.
\end{hint}

\begin{hint}{pb:shifted}
(a) $p\equiv1\pmod q$ has density $\frac1{q-1}$ (Tome~5, Theorem~5.8), and conditions for distinct $q$ combine by the Chinese remainder theorem. (c) At most one prime
$q>\sqrt x$ divides $p-1$; split at $x^\varepsilon$. (e) At most one large prime factor.
\end{hint}

\begin{hint}{pb:romanoff}
(a) $\pi(\frac x2)\cdot\lfloor\log_2\frac x2\rfloor$ pairs give $n\le x$. Cauchy--Schwarz. (b) Off the diagonal, $p'=p+2^l(2^m-1)$. (c) $1+g(q)=\bigl(\frac q{q-1}\bigr)^2$.
$e\mid2^m-1$ exactly when $\operatorname{ord}_e(2)\mid m$. (d) $P_J$ has fewer than $J^2$ prime factors; Mertens. (e) The orders of $2$ modulo $3,5,7,13,17,241$ are
$2,4,3,12,8,24$.
\end{hint}

\begin{hint}{pb:dedekind}
(a) The sawtooth on $\Z/k\Z$ has discrete Fourier coefficients $\frac i{2k}\cot\frac{\pi j}k$; Parseval. (b) Inside the rectangle the poles are at $z=1$ (triple) and at the
points $\frac jh$, $\frac jk$ with $0<j<h$, $0<j<k$; as $Y\to\infty$, the cotangents tend to $\mp i$ on the horizontal sides. (c) For $|s|\le s(1,k)$, Cauchy--Schwarz. (d) Reciprocity
along Euclid's algorithm.
\end{hint}

\section{Solutions}

\begin{solution}{pb:jacobi}
\idea Jacobi's identity squares to a double sum, and a change of variables turns that double sum into a combination of two theta series and their derivatives. Dividing
by the theta series turns the derivatives into logarithmic derivatives of products, which are Lambert series: divisor sums appear by themselves.

\writeup (a) Put $z=-1$ in the triple product: $\Theta(-x)=\sum(-1)^mx^{m^2}=\prod(1-x^{2n})(1-x^{2n-1})^2$. Since $\prod(1-x^{2n-1})=\frac{\prod(1-x^n)}{\prod(1-x^{2n})}=\frac1{\prod(1+x^n)}$
and $1-x^{2n}=(1-x^n)(1+x^n)$, the product is $\prod(1-x^n)(1+x^n)\cdot\prod(1+x^n)^{-2}=\prod\frac{1-x^n}{1+x^n}$.

(b) Replacing $k$ by $-1-k$ leaves $(-1)^k(2k+1)$ and $\frac{k(k+1)}2$ unchanged, so $E^3=\frac12\sum_{k\in\Z}(-1)^k(2k+1)x^{k(k+1)/2}$ and
\[
E^6=\frac14\sum_{k,l\in\Z}(-1)^{k+l}(2k+1)(2l+1)x^{(k^2+k+l^2+l)/2}.
\]
If $k+l$ is even, put $k=r+s$, $l=r-s$: then $(2k+1)(2l+1)=(2r+1)^2-4s^2$ and $\frac{k^2+k+l^2+l}2=r^2+r+s^2$. If $k+l$ is odd, put $k=s+r$, $l=s-r-1$: then the sign is $-1$,
$(2k+1)(2l+1)=4s^2-(2r+1)^2$ and the exponent is again $r^2+r+s^2$. Both substitutions are bijections onto $\Z^2$, so each $(r,s)$ receives $2\bigl((2r+1)^2-4s^2\bigr)$, and
$E^6=\frac12\sum_{r,s}\bigl((2r+1)^2-4s^2\bigr)x^{r^2+r+s^2}$.

(c) Since $(2r+1)^2=4(r^2+r)+1$, the double sum of (b) is $\sum_r\bigl(4(r^2+r)+1\bigr)x^{r^2+r}\cdot\sum_sx^{s^2}-4\sum_rx^{r^2+r}\cdot\sum_ss^2x^{s^2}=(4DB+B)\Theta-4B\,D\Theta$.
The triple product with $z=1$ gives the product for $\Theta$. With $z=x$ it gives $\sum_mx^{m^2+m}=\prod(1-x^{2n})(1+x^{2n})(1+x^{2n-2})$, and the factor $n=1$ of the last product is
$2$, so $B=2\prod(1-x^{2n})(1+x^{2n})^2$.

(d) Divide (c) by $\Theta B$: $\frac{2E^6}{\Theta B}=1+4\frac{DB}B-4\frac{D\Theta}\Theta$. By the products,
$\Theta B=2\prod(1-x^{2n})^2(1+x^n)^2$, so $\frac{2E^6}{\Theta B}=\prod\frac{(1-x^n)^6}{(1-x^n)^2(1+x^n)^4}=\Theta(-x)^4$ by (a). Logarithmic derivatives of the products give
\[
\frac{DB}B-\frac{D\Theta}\Theta=\sum_{n\ge1}\Bigl(\frac{4nx^{2n}}{1+x^{2n}}-\frac{2(2n-1)x^{2n-1}}{1+x^{2n-1}}\Bigr),
\]
the terms $-\frac{2nx^{2n}}{1-x^{2n}}$ cancelling. This is the formula for $\Theta(-x)^4$. Now replace $x$ by $-x$:
\[
\Theta(x)^4=1+8\sum_{m\text{ odd}}\frac{mx^m}{1-x^m}+8\sum_{m\text{ even}}\frac{mx^m}{1+x^m}.
\]
Expanding, $r_4(N)=8\sum_{m\mid N}m\,\varepsilon(m)$ with $\varepsilon=1$ for odd $m$ and $\varepsilon=(-1)^{N/m-1}$ for even $m$. For odd $N$ this is $8\sigma(N)$. For
$N=2^au$, $a\ge1$, $u$ odd, the even divisors $2^bv$ ($v\mid u$) have sign $-1$ for $b<a$ and $+1$ for $b=a$, so the total is
$8\sigma(u)\bigl(1+2^a-(2+4+\dots+2^{a-1})\bigr)=8\cdot3\sigma(u)$, which is $8\sum_{d\mid N,4\nmid d}d$ again.

(e) $r_4(n)\ge8>0$, so every $n\ge1$ is a sum of four squares. $r_4(p)=8(p+1)$, and $r_4(2^k)=8(1+2)=24$ for $k\ge1$.

(f) Let $R=\sqrt x$. The points $v\in\Z^4$ with $|v|\le R$ number $1+\sum_{n\le x}r_4(n)$, and the unit cubes $v+[0,1[^4$ cover the ball of radius $R-2$ and lie in the
ball of radius $R+2$. The volumes are $\frac{\pi^2}2(R\pm2)^4=\frac{\pi^2}2x^2+O(x^{3/2})$. On the other hand, by (d),
$\sum_{n\le x}r_4(n)=8\sum_{m\le x}\sum_{d\le x/m,\,4\nmid d}d$, and $\sum_{d\le y,\,4\nmid d}d=\frac{y^2}2-4\cdot\frac12\bigl(\frac y4\bigr)^2+O(y)=\frac38y^2+O(y)$. Hence
\[
\sum_{n\le x}r_4(n)=3x^2\sum_{m\le x}\frac1{m^2}+O\Bigl(x\sum_{m\le x}\frac1m\Bigr)=3\zeta(2)x^2+O(x\log x).
\]
Comparing the two leading terms, $3\zeta(2)=\frac{\pi^2}2$.

\tools{Jacobi's triple product and identity (Problem~\ref{pb:ramanujan5}); Lambert series; lattice points and volumes (Tome~0); $\sum_{m\le x}\frac1{m^2}=\zeta(2)+O(\frac1x)$.}

\leads Jacobi found the formula in the \emph{Fundamenta Nova} of 1829 with elliptic functions; the elementary route through $E^6$ is Hirschhorn's (1987). The two counts in (f)
give a proof of $\zeta(2)=\frac{\pi^2}6$ by counting points in a four-dimensional ball, with the volume $\frac{\pi^2}2R^4$ carrying the $\pi$. In modular language, $\Theta(\tau)^4$ is a modular form of weight $2$ for $\Gamma_0(4)$, a two-dimensional space spanned by Eisenstein
series, which is why a pure divisor sum comes out.
\end{solution}

\begin{solution}{pb:densities}
\idea Modulo a high power of $p$ a solution can be lifted coordinate by coordinate, as in Newton's method, provided one coordinate is a unit. So the number of solutions
modulo $p^k$ grows exactly like $p^{3k}$, and the constant of proportionality is read off modulo $p$ (or modulo $8$ when $p=2$).

\writeup Write $Q(x)=x_1^2+x_2^2+x_3^2+x_4^2$. (a) Let $x$ be a primitive solution modulo $p^k$, $k\ge1$, say $p\nmid x_1$, and $Q(x)=n+p^kc$. Its $p^4$ lifts
$x+p^kt$ modulo $p^{k+1}$ satisfy $Q(x+p^kt)\equiv n+p^k(c+2\sum x_it_i)\pmod{p^{k+1}}$, since $2k\ge k+1$. The condition $c+2\sum x_it_i\equiv0\pmod p$ is a linear equation in $t$
with the unit coefficient $2x_1$, so exactly $p^3$ lifts are solutions; every primitive solution modulo $p^{k+1}$ arises once in this way. A non-primitive solution is $x=py$,
with $y$ determined modulo $p^{k-1}$, and $Q(x)=p^2Q(y)$. If $p^2\nmid n$ there is none ($k\ge2$). If $p^2\mid n$, the condition is $Q(y)\equiv\frac n{p^2}\pmod{p^{k-2}}$, which
depends on $y$ modulo $p^{k-2}$; each such class has $p^4$ lifts modulo $p^{k-1}$, so there are $p^4N_{p^{k-2}}(\frac n{p^2})$ of them.

(b) By (a), $p^{-3k}N^*_{p^k}(n)=p^{-3}N^*_p(n)=:D^*$ for all $k\ge1$. By Problem~\ref{pb:sumsquares}, $N_p(n)=p^3-p$ if $p\nmid n$, and $N_p(n)=p^3+p^2-p$ if $p\mid n$; removing the
zero vector in the second case, $D^*=1-p^{-2}$ if $a=0$, and $D^*=p^{-3}(p^3+p^2-p-1)=(1-p^{-2})(1+p^{-1})$ if $a\ge1$. For $a\le1$ all solutions modulo $p^k$, $k\ge2$, are
primitive, so $\delta_p(n)=D^*$. For $a\ge2$, (a) gives $p^{-3k}N_{p^k}(n)=D^*+p^{-2}\,p^{-3(k-2)}N_{p^{k-2}}(\frac n{p^2})$, so by induction on $a$ the limit exists and
$\delta_p(n)=(1-p^{-2})(1+p^{-1})+p^{-2}(1-p^{-2})(1+\dots+p^{-(a-2)})=(1-p^{-2})(1+p^{-1}+\dots+p^{-a})$.

(c) Let $x$ be primitive modulo $2^k$, $k\ge3$, with $x_i$ odd. All $16$ lifts $x+2^kt$ have $Q\equiv Q(x)+2^{k+1}\sum x_it_i\equiv Q(x)\pmod{2^{k+1}}$, so either all or none are
solutions modulo $2^{k+1}$. Let $x'=x+2^{k-1}e_i$, where $i$ is the first odd coordinate: $Q(x')=Q(x)+2^kx_i+2^{2k-2}\equiv Q(x)+2^k\pmod{2^{k+1}}$ as $2k-2\ge k+1$. So
$x\mapsto x'$ is an involution without fixed points on the primitive solutions modulo $2^k$, and in each pair exactly one element has its $16$ lifts solutions. Hence
$N^*_{2^{k+1}}=\frac{16}2N^*_{2^k}$. Modulo $8$, odd squares are $1$ and even squares are $0$ (for $0,4$) or $4$ (for $2,6$). If $j$ coordinates are odd and $m$ of the even ones are
$\equiv2\pmod4$, then $Q\equiv j+4m\pmod8$. Counting, $N^*_8(n)=4\cdot4\cdot32=512$ for odd $n$ ($j\equiv n\pmod4$), $6\cdot16\cdot8=768$ for $n\equiv2\pmod4$ ($j=2$), $256$ for
$n\equiv4\pmod8$ ($j=4$), and $0$ for $8\mid n$. So $D^*=2^{-9}N^*_8(n)=1,\frac32,\frac12,0$ in these cases, and as in (b), $\delta_2(n)=D^*+\frac14\delta_2(\frac n4)$ when $4\mid n$. This
gives $\delta_2=1$ for $a=0$, $\frac32$ for $a=1$, $\frac12+\frac14=\frac34$ for $a=2$, and $\delta_2(n)=\frac14\delta_2(\frac n4)$ for $a\ge3$; by induction $\delta_2(n)=3\cdot2^{-a}$ for $a\ge1$.

(d) For odd $p\nmid n$, $\delta_p=1-p^{-2}$, and $\prod_{p>2}(1-p^{-2})=\frac{6/\pi^2}{3/4}=\frac8{\pi^2}$. Let $n=2^au$ with $u$ odd. Since $(1+p^{-1}+\dots+p^{-a_p})=\frac{\sigma(p^{a_p})}{p^{a_p}}$,
\[
\pi^2n\prod_p\delta_p(n)=\pi^2\cdot2^au\cdot\delta_2(n)\cdot\frac8{\pi^2}\cdot\frac{\sigma(u)}u=8\cdot2^a\delta_2(n)\,\sigma(u),
\]
which is $8\sigma(u)$ for $a=0$ and $24\sigma(u)$ for $a\ge1$: exactly $r_4(n)$ by Problem~\ref{pb:jacobi}(d). The product converges absolutely because all but finitely many factors are
$1-p^{-2}$. The ball $Q(x)\le t$ has volume $\frac{\pi^2}2t^2$, whose derivative at $t=n$ is $\pi^2n$: the real solutions of $Q(x)=n$ have ``density'' $\pi^2n$, just as
$p^{-3k}N_{p^k}(n)$ is the density of solutions modulo $p^k$.

\tools{Hensel's lifting (Tome~1); $N_p(n)$ from Problem~\ref{pb:sumsquares}; $\zeta(2)=\frac{\pi^2}6$; Problem~\ref{pb:jacobi}.}

\leads ``Number of solutions $=$ density at infinity $\times$ product of local densities'' is the Hardy--Littlewood circle method in its simplest form. For four squares the
formula is exact; for a positive definite form $Q$ in general, Siegel's mass formula (1935) says that it holds for a weighted average over the forms of the same genus, and
the circle method proves it asymptotically for five or more variables. Problem~\ref{pb:jacobi} computed the left side and this problem the right side, independently.
\end{solution}

\begin{solution}{pb:twosquares}
\idea Squaring $F$ produces $\zeta(s)L(s,\chi_4)$, so $F$ behaves like $(s-1)^{-1/2}$: a ``half pole''. Rankin's trick turns this into the size $\sqrt{\log y}$ of the logarithmic
sum $S(y)$, and two elementary arguments, one through prime powers and one through primes $p\equiv1\pmod4$, transfer it to $B(x)$.

\writeup (a) By Tome~1, Theorem~7.3, $b$ is multiplicative, with $b(2^k)=1$, $b(p^k)=1$ for $p\equiv1\pmod4$, and $b(p^k)=[k\text{ even}]$ for $p\equiv3\pmod4$. So $F(s)$ has Euler
factors $(1-2^{-s})^{-1}$, $(1-p^{-s})^{-1}$ and $(1-p^{-2s})^{-1}$. The factors of $\zeta(s)L(s,\chi_4)$ are $(1-2^{-s})^{-1}$, $(1-p^{-s})^{-2}$ and $(1-p^{-s})^{-1}(1+p^{-s})^{-1}=(1-p^{-2s})^{-1}$.
Comparing squares gives the identity. As $\sigma\to1^+$, $\zeta(\sigma)\sim\frac1{\sigma-1}$, $L(1,\chi_4)=\frac\pi4$, $\frac1{1-2^{-1}}=2$, and the last product converges at $s=1$, so
$F(\sigma)^2\sim\frac{C^2}{\sigma-1}$.

(b) For $n\le y$, $n^{-\delta}\ge e^{-1}$, so $\frac1n\le e\,n^{-1-\delta}$ and $S(y)\le eF(1+\delta)$. For $n>y$, $n^{-1-\delta}\le y^{-\delta/2}n^{-1-\delta/2}=e^{-1/2}n^{-1-\delta/2}$, so
$S(y)\ge\sum_{n\le y}b(n)n^{-1-\delta}\ge F(1+\delta)-e^{-1/2}F(1+\frac\delta2)$. By (a), $F(1+\delta)\sim C\sqrt{\log y}$ and $F(1+\frac\delta2)\sim C\sqrt{2\log y}$, so
$C\bigl(1-\sqrt2e^{-1/2}+o(1)\bigr)\sqrt{\log y}\le S(y)\le C(e+o(1))\sqrt{\log y}$, and $1-\sqrt2e^{-1/2}=0.142\ldots>0$.

(c) For $p\le\sqrt y$ there are at most $\frac{\log y}{\log p}$ values of $k$, so $\sum_{p^k\le y}k\log p\le\sum_{p\le\sqrt y}\frac{\log^2y}{\log p}+\sum_{p\le y}\log p=O(\sqrt y\log y)+O(y)$
by Chebyshev (Tome~5, Proposition~1.2). If $b(n)=1$ and $p^k\parallel n$, then $m=\frac n{p^k}$ has $b(m)=1$. Writing $\log n=\sum_{p^k\parallel n}k\log p$ and counting $n$ by $m$,
\[
\sum_{n\le x}b(n)\log n\le\sum_{m\le x}b(m)\sum_{p^k\le x/m}k\log p\le c\sum_{m\le x}b(m)\frac xm=c\,x\,S(x).
\]
With $L=\log x$, $\log\frac xL\,\bigl(B(x)-B(\frac xL)\bigr)\le\sum_{n\le x}b(n)\log n=O(x\sqrt{\log x})$ by (b), so $B(x)\le\frac x{\log x}+O\bigl(\frac{x\sqrt{\log x}}{\log x}\bigr)=O\bigl(\frac x{\sqrt{\log x}}\bigr)$.

(d) If $b(m)=1$ and $p\equiv1\pmod4$, then $b(mp)=1$, and conversely $b(n)=1$, $p\mid n$ give $b(\frac np)=1$. So pairs $(m,p)$ with $mp\le x$ correspond to pairs $(n,p)$ with $p\mid n$, and
\[
\sum_{n\le x}b(n)\log n\ge\sum_{n\le x}b(n)\sum_{p\mid n,\,p\equiv1\,(4)}\log p=\sum_{m\le x}b(m)\,\theta\bigl(\tfrac xm;4,1\bigr)\ge\sum_{m\le x/\log^3x}b(m)\,\theta\bigl(\tfrac xm;4,1\bigr).
\]
For $m\le\frac x{\log^3x}$, $\frac xm\ge\log^3x\to\infty$, and $\theta(y;4,1)\sim\frac y2$ (Tome~5, Theorem~5.8), so the right side is at least $(\frac12-o(1))\,x\,S(\frac x{\log^3x})$, which
is $\ge c\,x\sqrt{\log x}$ by (b). Since $\sum_{n\le x}b(n)\log n\le B(x)\log x$, $B(x)\ge\frac{c\,x}{\sqrt{\log x}}$.

(e) By partial summation, for $\sigma>1$, $F(\sigma)=\sigma\int_1^\infty B(x)x^{-\sigma-1}\dd x$. If $B(x)=(K+o(1))\frac x{\sqrt{\log x}}$, the part $x\le X$ is bounded for each fixed $X$, and
with $x=e^u$,
\[
\int_X^\infty\frac{x^{-\sigma}}{\sqrt{\log x}}\dd x=\int_{\log X}^\infty\frac{e^{-(\sigma-1)u}}{\sqrt u}\dd u=\frac{\Gamma(\frac12)+o(1)}{\sqrt{\sigma-1}}\qquad(\sigma\to1^+).
\]
So $F(\sigma)\sim K\sqrt\pi(\sigma-1)^{-1/2}$, and (a) forces $K=\frac C{\sqrt\pi}$. Numerically $\prod_{p\equiv3\,(4)}(1-p^{-2})^{-1}=1.168076\ldots$, $C=1.35455\ldots$ and $K=0.764224\ldots$
The notebook's $0.7922$ is $1.037K$.

\tools{Tome~1, Theorem~7.3; Tome~5, Proposition~1.2 and Theorem~5.8; $L(1,\chi_4)=\frac\pi4$ (Leibniz); Rankin's trick (Problem~\ref{pb:rankin}); $\Gamma(\frac12)=\sqrt\pi$.}

\leads Landau proved in 1908 that $B(x)\sim\frac{Kx}{\sqrt{\log x}}$ indeed, with complex analysis: $F(s)$ is $\sqrt{\zeta(s)}$ times a function holomorphic near $s=1$, and he integrated
around the branch cut. Tome~7 proves the logarithmic version $\sum_{n\le x}\frac{b(n)}n\sim2K\sqrt{\log x}$ with Karamata's Tauberian theorem. The slow convergence seen in the
data comes from a secondary term: $B(x)=\frac{Kx}{\sqrt{\log x}}\bigl(1+\frac{c}{\log x}+\cdots\bigr)$, and at $x=10^8$ the correction $\frac c{\log x}$ is still $3.7\%$.
\end{solution}

\begin{solution}{pb:ramanujanE2}
\idea $E_2$ is almost a modular form: its transformation law has an extra term $\frac{6\tau}{\pi i}$, inherited from the $\sqrt{-i\tau}$ in Dedekind's law. Differentiating a modular
form also spoils modularity by a term of the same kind, and the two defects cancel in $DE_2-\frac{E_2^2}{12}$. Since the spaces $M_4$, $M_6$, $M_8$ are one-dimensional, each corrected
derivative is a multiple of a known form.

\writeup (a) $\log\eta=\frac{2\pi i\tau}{24}+\sum\log(1-q^n)$, so $D\log\eta=\frac1{24}-\sum_n\frac{nq^n}{1-q^n}=\frac1{24}-\sum_N\sigma(N)q^N=\frac{E_2}{24}$.

(b) Apply $D$ to $\log\eta(-\frac1\tau)=\log\eta(\tau)+\frac12\log(-i\tau)$. By the chain rule the left side gives $\frac1{\tau^2}\frac{E_2}{24}(-\frac1\tau)$, the right side
$\frac{E_2(\tau)}{24}+\frac1{2\pi i}\frac1{2\tau}$. Multiplying by $24\tau^2$: $E_2(-\frac1\tau)=\tau^2E_2(\tau)+\frac{12\tau}{2\pi i}=\tau^2E_2(\tau)+\frac{6\tau}{\pi i}$.

(c) Write $c=\frac{12}{2\pi i}$, so that $E_2(-\frac1\tau)=\tau^2E_2+c\tau$. Applying $D$ and the chain rule, $(DE_2)(-\frac1\tau)=\tau^4DE_2+\frac{2\tau^3E_2+c\tau^2}{2\pi i}$, and
$E_2(-\frac1\tau)^2=\tau^4E_2^2+2c\tau^3E_2+c^2\tau^2$. In $f(-\frac1\tau)$ the $\tau^3E_2$ terms are $\frac2{2\pi i}-\frac{2c}{12}=0$ and the $\tau^2$ terms are $c\bigl(\frac1{2\pi i}-\frac c{12}\bigr)=0$,
so $f(-\frac1\tau)=\tau^4f(\tau)$. Also $f(\tau+1)=f(\tau)$, and $f=-\frac1{12}+O(q)$ is holomorphic at $i\infty$. So $f\in M_4=\C E_4$ (Tome~5, Theorem~12.9), and $f=-\frac{E_4}{12}$.
In the same way, for a form $g$ of weight $k$, $(Dg)(-\frac1\tau)=\tau^{k+2}Dg+\frac{k\tau^{k+1}g}{2\pi i}$, so $Dg-\frac k{12}E_2g$ has weight $k+2$. For $g=E_4$ it lies in $M_6=\C E_6$
with constant term $-\frac13$; for $g=E_6$ it lies in $M_8=\C E_4^2$ with constant term $-\frac12$. These are Ramanujan's equations.

(d) The coefficient of $q^n$ in $12DE_2$ is $-288n\sigma(n)$; in $E_2^2$ it is
\[
-48\sigma(n)+576\sum_{m=1}^{n-1}\sigma(m)\sigma(n-m);
\]
in $E_4$ it is $240\sigma_3(n)$. Solving,
$576\sum\sigma(m)\sigma(n-m)=240\sigma_3(n)+(48-288n)\sigma(n)$, which is the formula after division by $48$.

(e) In $3DE_4=E_2E_4-E_6$, the coefficients of $q^n$ give $720n\sigma_3(n)=240\sigma_3(n)-24\sigma(n)-5760\sum\sigma(m)\sigma_3(n-m)+504\sigma_5(n)$; divide by $24$.

(f) With $\sigma(p)=p+1$ and $\sigma_3(p)=p^3+1$, (d) gives $\frac{5p^3+5+(1-6p)(p+1)}{12}=\frac{5p^3-6p^2-5p+6}{12}=\frac{(p-1)(5p^2-p-6)}{12}$. For $p=7$: $\sigma(1),\dots,\sigma(6)=1,3,4,7,6,12$ and
$1\cdot12+3\cdot6+4\cdot7+7\cdot4+6\cdot3+12\cdot1=116=\frac{6\cdot232}{12}$.

\tools{Tome~5, Theorem~12.7 (Dedekind's $\eta(-\frac1\tau)=\sqrt{-i\tau}\,\eta(\tau)$) and Theorem~12.9 (dimensions of $M_k$, so $E_8=E_4^2$).}

\leads Ramanujan published these equations in 1916, in the paper where he introduced $\tau(n)$. They say that the ring $\C[E_2,E_4,E_6]$ is closed under $D$; it is the ring of
\emph{quasimodular forms}. Every convolution sum $\sum\sigma_a(m)\sigma_b(n-m)$ with $a$ and $b$ odd can be evaluated this way, except that from weight $12$ on the cusp form $\Delta$
enters and its coefficients $\tau(n)$ appear in the answer.
\end{solution}

\begin{solution}{pb:landauzeros}
\idea Integrate $x^s\frac{\zeta'}\zeta(s)$ around a tall rectangle. The residues are the $x^\rho$. On the right edge the Dirichlet series of $\frac{\zeta'}\zeta$ is integrated term by term,
and only the term $n=x$ does not oscillate: it grows like $T$. Everything else is smaller, so the zeros, summed with the weights $x^\rho$, ``detect'' whether $x$ is a prime power.

\writeup Fix $x>1$ and $T\ge3$, and choose $T'\in[T,T+1]$ as in Tome~5, Lemma~7.3, so that $|\frac{\zeta'}\zeta(\sigma+iT')|\le C\log^2T$ for $-1\le\sigma\le3$. No zero has $0<\gamma\le1$.

(a) On $\operatorname{Re}s=2$ the series $\frac{\zeta'}\zeta(s)=-\sum\Lambda(n)n^{-s}$ converges absolutely, so
\[
\frac1{2\pi}\int_1^Tx^{2+it}\frac{\zeta'}\zeta(2+it)\dd t=-\frac1{2\pi}\sum_n\Lambda(n)\Bigl(\frac xn\Bigr)^2\int_1^T\Bigl(\frac xn\Bigr)^{it}\dd t.
\]
The term $n=x$ (if $x$ is an integer) gives $-\frac{T-1}{2\pi}\Lambda(x)$. For $n\ne x$ the integral is at most $\frac2{|\log(x/n)|}$, and $\sum_{n\ne x}\Lambda(n)\frac{x^2}{n^2|\log(x/n)|}$
converges, since $|\log\frac xn|$ is bounded below for $n\ne x$ and $\sum\frac{\Lambda(n)}{n^2}<\infty$. The total is $-\frac T{2\pi}\Lambda(x)+O_x(1)$.

(b) Take the logarithmic derivative of $\xi(s)=\xi(1-s)$ with $\xi(s)=\frac12s(s-1)\pi^{-s/2}\Gamma(\frac s2)\zeta(s)$:
\[
\frac1s+\frac1{s-1}-\frac{\log\pi}2+\frac12\frac{\Gamma'}\Gamma\Bigl(\frac s2\Bigr)+\frac{\zeta'}\zeta(s)=\frac1s+\frac1{s-1}+\frac{\log\pi}2-\frac12\frac{\Gamma'}\Gamma\Bigl(\frac{1-s}2\Bigr)-\frac{\zeta'}\zeta(1-s),
\]
which is the stated formula. On $s=-1+it$, $\frac{\zeta'}\zeta(1-s)=\frac{\zeta'}\zeta(2-it)$ is again an absolutely convergent series, and integrating $x^{-1+it}\Lambda(n)n^{-2+it}$ gives
$O(\frac1{\log(xn)})$, summable in $n$. The constant $\log\pi$ gives $O(\frac1{\log x})$. The function $\varphi(t)=-\frac12\frac{\Gamma'}\Gamma(\frac{-1+it}2)-\frac12\frac{\Gamma'}\Gamma(\frac{2-it}2)$ satisfies
$\varphi(t)=O(\log t)$ and $\varphi'(t)=O(\frac1t)$ for $t\ge1$, by Stirling's formula for $\frac{\Gamma'}\Gamma$ (Tome~5, Theorem~6.1) and the reflection formula to move
$\frac{-1+it}2$ to the right half-plane. Integrating by parts,
$\int_1^{T}x^{it}\varphi(t)\dd t=\bigl[\frac{x^{it}}{i\log x}\varphi\bigr]_1^{T}-\int_1^{T}\frac{x^{it}}{i\log x}\varphi'(t)\dd t=O_x(\log T)$.

(c) The rectangle with vertices $2+i$, $2+iT'$, $-1+iT'$, $-1+i$ contains the zeros with $1<\gamma<T'$ and no other pole of $x^s\frac{\zeta'}\zeta$; the residue at $\rho$ is
$x^\rho$ times its multiplicity. So $\sum_{0<\gamma<T'}x^\rho$ is $\frac1{2\pi i}$ times the integral around the boundary. The right edge gives $-\frac{T'}{2\pi}\Lambda(x)+O_x(1)$ by (a), the
left edge $O_x(\log T)$ by (b), the bottom edge $O_x(1)$, and the top edge at most $\frac1{2\pi}\cdot3x^2\cdot C\log^2T$. Finally $T'-T\le1$, and there are $O(\log T)$ zeros with
$T<\gamma\le T'$ (Tome~5, Corollary~6.10), each with $|x^\rho|\le x$. This gives Landau's formula.

(d) Under the Riemann hypothesis $x^\rho=\sqrt x\,e^{i\gamma\log x}$. Take real parts in (c) and divide by $\sqrt x$.

(e) By Weyl's criterion (Tome~4, Theorem~6.13) it is enough to show that $\frac1N\sum_{n\le N}e^{2\pi ik\alpha\gamma_n}\to0$ for every integer $k\ne0$; replacing $k$ by $-k$ conjugates the sum,
so let $k\alpha>0$, and $y=e^{2\pi k\alpha}>1$. Then $e^{2\pi ik\alpha\gamma}=y^{i\gamma}=y^{-1/2}y^\rho$, so by (c)
$\sum_{0<\gamma\le T}e^{2\pi ik\alpha\gamma}=-\frac T{2\pi}\frac{\Lambda(y)}{\sqrt y}+O_y(\log^2T)=O(T)$, while the number of terms is $N(T)\sim\frac T{2\pi}\log T$ (Tome~5, Theorem~6.9).
The ratio tends to $0$, also along $T=\gamma_N$.

\tools{Tome~5, Proposition~6.2, Theorem~6.1, Lemma~7.3, Theorem~6.9 and Corollary~6.10; Weyl's criterion (Tome~4, Theorem~6.13).}

\leads Landau proved his formula in 1911; Gonek (1993) made it uniform in $x$, which is what one needs to use it as a tool. Part (e) is Rademacher's theorem (1956); Hlawka (1975)
removed the Riemann hypothesis. The spikes of Figure~\ref{fig:landau} are the formula with $2000$ zeros: at a prime power $p^k$ the value is within $3\%$ of $\log p$, and between the
spikes it is a small ripple.
\end{solution}

\begin{solution}{pb:hearzeros}
\idea The explicit formula is the residue theorem applied to $H\frac{\zeta'}\zeta$, where $H$ is a smooth bump on the critical line that is still entire. The zeros come out as
residues, the primes from the Dirichlet series on the right edge, and the functional equation folds the left edge back onto the right. Choosing $H$ concentrated near height
$T$ turns the formula into a zero detector built from primes.

\writeup (a) $h(w)=e^{-(w-T)^2/2s^2}+e^{-(w+T)^2/2s^2}$ is entire and even in $w$, so $H$ is entire and $H(1-z)=h(i(z-\frac12))=h(-i(z-\frac12))=H(z)$; and $H(\frac12+it)=h(t)$.
For $z=\sigma+it$, $w=t-i(\sigma-\frac12)$, and $|e^{-(w\mp T)^2/2s^2}|=e^{-((t\mp T)^2-(\sigma-\frac12)^2)/2s^2}$, a Gaussian in $t$ uniformly for bounded $\sigma$. The value of $g$ is the
Fourier transform of a Gaussian: $\frac1{2\pi}\int e^{-(t-T)^2/2s^2}e^{-iut}\dd t=\frac s{\sqrt{2\pi}}e^{-iuT}e^{-s^2u^2/2}$, plus the same with $-T$.

(b) Integrate $H\frac{\zeta'}\zeta$ around the rectangle $-1\le\sigma\le2$, $|t|\le T'$, with good heights $T'\to\infty$ (Tome~5, Lemma~7.3): by (a) the horizontal sides tend to $0$,
and the residues are $H(\rho)$ at the zeros and $-H(1)$ at the pole of $\zeta$. So
$\sum_\rho H(\rho)-H(1)=\frac1{2\pi i}\int_{(2)}H\frac{\zeta'}\zeta-\frac1{2\pi i}\int_{(-1)}H\frac{\zeta'}\zeta$, both lines upward.

\emph{Right edge.} Expanding $\frac{\zeta'}\zeta$ and moving each term to the line $\frac12$ (the integrand $Hn^{-z}$ is entire and decays),
\[
\frac1{2\pi i}\int_{(2)}H(z)n^{-z}\dd z=\frac1{2\pi}\int h(t)n^{-1/2-it}\dd t=n^{-1/2}g(\log n).
\]
The right edge is $-\sum\Lambda(n)n^{-1/2}g(\log n)$.

\emph{Left edge.} Insert the formula of Problem~\ref{pb:landauzeros}(b). The term $-\frac{\zeta'}\zeta(1-z)$, after $z\mapsto1-z$ and $H(1-z)=H(z)$, gives the right edge
once more. The constant $\log\pi$ gives $-\log\pi\cdot\frac1{2\pi}\int h=-g(0)\log\pi$. The remaining term is $\frac1{2\pi i}\int_{(-1)}H(z)\cdot\frac12\bigl(\frac{\Gamma'}\Gamma(\frac z2)+\frac{\Gamma'}\Gamma(\frac{1-z}2)\bigr)\dd z$.
Moving it to the line $\frac12$ crosses one pole, at $z=0$, where $\frac12\frac{\Gamma'}\Gamma(\frac z2)\sim-\frac1z$; so it equals $H(0)$ plus
$\frac1{2\pi}\int h(t)\operatorname{Re}\frac{\Gamma'}\Gamma(\frac14+\frac{it}2)\dd t$, because $\frac{1-z}2=\overline{\frac z2}$ on that line. Collecting terms gives the formula.
It converges absolutely because $H(\rho)$ decays faster than any power of $|\gamma|$ while there are $O(T\log T)$ zeros up to height $T$.

(c) Under the Riemann hypothesis $H(\rho)=h(\gamma)$, and $\sum_\rho h(\gamma)=2\sum_{\gamma>0}h(\gamma)=2Z(T)+O(e^{-cT})$, since $e^{-(\gamma+T)^2/2s^2}\le e^{-\gamma^2/2s^2}e^{-14T/s^2}$.
Also $|H(0)|=|H(1)|=|h(\frac i2)|\le2e^{(1/4-T^2)/2s^2}$. In the integral, $h$ is concentrated near $\pm T$, where
$\operatorname{Re}\frac{\Gamma'}\Gamma(\frac14+\frac{it}2)=\log\frac{|t|}2+O(t^{-2})$ (Stirling). Expanding $\log\frac t2=\log\frac T2+\frac{t-T}T+O(\frac{(t-T)^2}{T^2})$ against the Gaussian,
the linear term integrates to $0$, so the integral is $\frac{2s}{\sqrt{2\pi}}\log\frac T2+O(\frac1{T^2})$. With $g(0)=\frac{2s}{\sqrt{2\pi}}$ the formula becomes
\[
2Z(T)=\frac{2s}{\sqrt{2\pi}}\log\frac T{2\pi}-2\sum_n\frac{\Lambda(n)}{\sqrt n}g(\log n)+O\Bigl(\frac1T\Bigr),
\]
and dividing by $2$ gives the claim.

(d) A zero $\gamma$ contributes a bump of height $1$ at $T=\gamma$, and the other zeros add little, so the peaks of Figure~\ref{fig:hearzeros} sit at $14.13$, $21.02$, $25.01$, \dots{}
The closest pair below $50$, $48.01$ and $49.77$, is $1.76=2.9s$ apart and still resolved. A smaller $s$ makes narrower bumps and separates closer zeros, but $g(\log n)$ decays only
like $e^{-s^2\log^2n/2}$, so the primes up to about $e^{c/s}$ are needed: halving the width squares the range of primes. This is the uncertainty principle: sharp in $t$ means wide
in $\log n$. The notebook checks that the right side agrees with $Z(T)$ computed from the zeros to within $0.003$ on $[5,50]$.

\tools{Problem~\ref{pb:landauzeros}(b); Tome~5, Lemma~7.3 and Theorem~6.1; Fourier transform of a Gaussian (Tome~4).}

\leads Guinand (1948) and Weil (1952) stated the explicit formula for general test functions; Weil showed that the Riemann hypothesis is equivalent to the positivity of a
quadratic form built from it. Zeros are located in practice with the Riemann--Siegel formula, which Siegel found in 1932 in Riemann's unpublished notes; the computation
of (d) is a different route, which uses nothing but the primes.
\end{solution}

\begin{solution}{pb:cramersign}
\idea In Cramér's model, $\pi_{\mathcal C}(x)-\operatorname{li}(x)$ is a sum of independent centred coin tosses, a random walk with no preference of sign. The real primes have a preference,
and its source is visible in (c): $\pi$ counts primes, while the explicit formula naturally counts prime powers with weights $\frac1k$. The squares of primes, which $\pi$ omits,
push $\pi(x)$ below $\operatorname{li}(x)$ by about $\frac{\sqrt x}{\log x}$, which is also the size of the honest fluctuations.

\writeup (a) $\mathbb E\,\pi_{\mathcal C}(x)=\sum_{3\le n\le x}\frac1{\log n}$ differs from $\int_3^x\frac{\dd t}{\log t}$ by at most $\frac1{\log3}$ (a decreasing function), so it is
$\operatorname{li}(x)+O(1)$. $\operatorname{Var}\pi_{\mathcal C}(x)=\sum\frac1{\log n}\bigl(1-\frac1{\log n}\bigr)\sim\frac x{\log x}$. The summands are bounded and the variance tends to infinity, so
Lindeberg's condition holds and $\frac{\pi_{\mathcal C}(x)-\mathbb E\pi_{\mathcal C}(x)}{\sqrt{\operatorname{Var}}}\to\mathcal N(0,1)$ (Tome~3, Theorem~6.10). As
$\frac{\operatorname{li}(x)-\mathbb E\pi_{\mathcal C}(x)}{\sqrt{\operatorname{Var}}}\to0$, $P(\pi_{\mathcal C}(x)>\operatorname{li}(x))\to P(Z>0)=\frac12$.

(b) Let $V_j=\operatorname{Var}\pi_{\mathcal C}(n_j)$ and let $D_j$ be the centred increment of $\pi_{\mathcal C}$ over $]n_{j-1},n_j]$. The $D_j$ are independent, $\operatorname{Var}D_j\sim V_j$ (as
$n_{j-1}=\sqrt{n_j}$), and by the central limit theorem $P(D_j>2\sqrt{V_j})\to P(Z>2)>0$. So $\sum_jP(D_j>2\sqrt{V_j})=\infty$, and by the second Borel--Cantelli lemma (Tome~3,
Theorem~4.11) $D_j>2\sqrt{V_j}$ for infinitely many $j$, almost surely; likewise $D_j<-2\sqrt{V_j}$. The centred value $Y_{j-1}$ at $n_{j-1}$ satisfies
$P(|Y_{j-1}|>\sqrt{V_j})\le\frac{V_{j-1}}{V_j}=O(n_j^{-1/2})$ by Chebyshev, a summable sequence, so by the first lemma $|Y_{j-1}|\le\sqrt{V_j}$ for all large $j$. Then
$\pi_{\mathcal C}(n_j)-\operatorname{li}(n_j)=Y_{j-1}+D_j+O(1)$ is $>\sqrt{V_j}-O(1)>0$ infinitely often and $<0$ infinitely often.

(c) We have
\[
\Pi(x)=\pi(x)+\frac12\pi(\sqrt x)+\sum_{3\le k\le\log_2x}\frac1k\pi(x^{1/k}),
\]
and the last sum is at most $x^{1/3}+x^{1/4}\log_2x=O(x^{1/3})$. Next, $\Pi(x)=\sum_{n\le x}\frac{\Lambda(n)}{\log n}$,
so by Stieltjes integration by parts, $\Pi(x)=\frac{\psi(x)}{\log x}+\int_2^x\frac{\psi(t)}{t\log^2t}\dd t$. In the same way
$\operatorname{li}(x)-\operatorname{li}(2)=\int_2^x\frac{\dd t}{\log t}=\frac x{\log x}-\frac2{\log2}+\int_2^x\frac{t}{t\log^2t}\dd t$. Subtracting gives the identity.

(d) By (c), $\pi(x)-\operatorname{li}(x)=\bigl(\Pi(x)-\operatorname{li}(x)\bigr)-\frac12\pi(\sqrt x)+O(x^{1/3})$, and $\frac12\pi(\sqrt x)\approx\frac{\sqrt x}{\log x}$. The term $\Pi-\operatorname{li}$ is dominated by
$\frac{\psi(x)-x}{\log x}$, which oscillates around $0$ (its $66\,112$ sign changes below $10^8$); the shift by $-\frac{\sqrt x}{\log x}$, about one unit on the scale of the data, is larger
than the oscillation in this range, so $\pi<\operatorname{li}$ throughout. Under the Riemann hypothesis, the explicit formula gives $\psi(x)-x=-\sqrt x\sum_\rho\frac{x^{i\gamma}}\rho+O(\log^2x)$,
so $\frac{\psi(x)-x}{\log x}$ has size $\frac{\sqrt x}{\log x}$ times an oscillating sum whose terms decay like $\frac1{|\gamma|}$. The integral in (c) is smaller: integrating $t^{\rho-1}$ against
$\frac1{\log^2t}$ divides each term by about $|\rho|\log x$, leaving $O(\frac{\sqrt x}{\log^2x})$. So both $\Pi-\operatorname{li}$ and $\pi-\operatorname{li}$ have fluctuations of size $\frac{\sqrt x}{\log x}$,
which is $543$ at $10^8$, while the coin model fluctuates by $\sqrt{x/\log x}=2330$: the real primes are far more regular than independent coins, because their errors are all
governed by the same zeros.

\tools{Tome~3, Theorems~4.11 and~6.10; Tome~5, Definition~10.1 and Theorem~7.5; Stieltjes integration by parts (Tome~5, Lemma~1.1).}

\leads The bias does not last forever: Littlewood proved in 1914 that $\pi(x)-\operatorname{li}(x)$ changes sign infinitely often, because $\sum_\rho\frac{x^{i\gamma}}\rho$ occasionally beats
the shift $-1$. But the first sign change is beyond $10^{19}$ (Büthe, 2018, verified $\pi(x)<\operatorname{li}(x)$ for $2\le x\le10^{19}$). Rubinstein and Sarnak (1994) showed, under the
Riemann hypothesis and a linear independence hypothesis on the $\gamma$, that the set of $x$ with $\pi(x)>\operatorname{li}(x)$ has logarithmic density $0.00000026\ldots$: Chebyshev's bias in
its sharpest form.
\end{solution}

\begin{solution}{pb:selberg}
\idea $\Lambda=\mu*\log$ sees one prime at a time; $\Lambda_2=\mu*\log^2$ sees primes and products of two primes, with total weight $2\log x$ on average. The point is that
$\sum\Lambda_2$ can be computed from Möbius sums alone, with an error $O(x)$, because the $\log^2$ is smooth enough for the error terms to cancel.

\writeup (a) $f\mapsto f\log$ is a derivation of the convolution ring: $(f*g)\log=(f\log)*g+f*(g\log)$, since $\log(dm)=\log d+\log m$. From $\log=1*\Lambda$ we get
$\log^2=\log*\Lambda+1*(\Lambda\log)$; convolving with $\mu$ and using $\mu*\log=\Lambda$, $\mu*1=\delta$ gives $\Lambda_2=\Lambda*\Lambda+\Lambda\log$. Both terms are $\ge0$, and both vanish when
$n$ has three distinct prime factors, since $\Lambda(d)\Lambda(m)\ne0$ forces $d,m$ to be prime powers.

(b) From $\sum_{m\le y}\frac1m=\log y+\gamma+O(\frac1y)$,
\[
\sum_{d\le x}\frac{\mu(d)}d\log\frac xd=\sum_{d\le x}\frac{\mu(d)}d\sum_{m\le x/d}\frac1m-\gamma\sum_{d\le x}\frac{\mu(d)}d+O(1),
\]
which is $\sum_{n\le x}\frac1n\sum_{d\mid n}\mu(d)+O(1)=O(1)$, using
$|\sum_{d\le x}\frac{\mu(d)}d|\le1$ (Tome~5, Exercise~1.6). For the second sum, the hyperbola method (Tome~5, Lemma~2.8) gives
$\sum_{m\le y}\frac{\tau(m)}m=\frac12\log^2y+2\gamma\log y+c+O(y^{-1/2}\log y)$. So $\log^2\frac xd=2\sum_{m\le x/d}\frac{\tau(m)}m-4\gamma\log\frac xd-2c+O\bigl((\frac dx)^{1/2}\log\frac xd\bigr)$. Multiply
by $\frac{\mu(d)}d$ and sum: since $\mu*\tau=1$, the first part is $2\sum_{n\le x}\frac1n=2\log x+O(1)$; the others are $O(1)$ by the first sum and Exercise~1.6, and the error
terms add up to $O\bigl(\sum_{d\le x}\frac1d(\frac dx)^{1/2}\log\frac xd\bigr)=O(1)$.

(c) $\sum_{n\le x}\Lambda_2(n)=\sum_{d\le x}\mu(d)\sum_{m\le x/d}\log^2m$, and $\sum_{m\le y}\log^2m=y(\log^2y-2\log y+2)+O(\log^2y+1)$ by comparison with the integral. So
\[
\sum_{n\le x}\Lambda_2(n)=x\sum_{d\le x}\frac{\mu(d)}d\Bigl(\log^2\frac xd-2\log\frac xd+2\Bigr)+O\Bigl(\sum_{d\le x}\bigl(\log^2\tfrac xd+1\bigr)\Bigr),
\]
which is $x\bigl(2\log x+O(1)\bigr)+O(x)$ by (b), Exercise~1.6, and $\sum_{d\le x}\log^2\frac xd=O(x)$.

(d) By (a),
\[
\sum_{n\le x}\Lambda_2(n)=\sum_{n\le x}\Lambda(n)\log n+\sum_{dm\le x}\Lambda(d)\Lambda(m).
\]
Prime powers $p^k$, $k\ge2$, contribute $O(\sqrt x\log^2x)$ to the first sum and at most
\[
\sum_{p^k,\,k\ge2}\log p\,\psi\Bigl(\frac x{p^k}\Bigr)=O\Bigl(x\sum_{p,\,k\ge2}\frac{\log p}{p^k}\Bigr)=O(x)
\]
to the second. And $\sum_{p\le x}\log^2p=\theta(x)\log x-\int_2^x\frac{\theta(t)}t\dd t=\theta(x)\log x+O(x)$.

(e) Write the formula as $\theta(x)\log x+\sum_{p\le x}\log p\,\theta(\frac xp)=2x\log x+O(x)$; by Chebyshev, $A<\infty$. Let $\varepsilon>0$; there is $X$ with $\theta(y)\le(A+\varepsilon)y$ for $y\ge X$,
and $\theta(\frac xp)=O_\varepsilon(1)$ for $p>\frac xX$. By Mertens (Tome~5, Theorem~1.3), $\sum_{p\le x}\log p\,\theta(\frac xp)\le(A+\varepsilon)x\log x+O_\varepsilon(x)$. Along a sequence with
$\theta(x)\le(a+\varepsilon)x$, the formula gives $2\le a+A+2\varepsilon+o(1)$. Symmetrically, $\theta(y)\ge(a-\varepsilon)y$ for $y\ge X$ gives
$\sum_p\log p\,\theta(\frac xp)\ge(a-\varepsilon)x\sum_{p\le x/X}\frac{\log p}p=(a-\varepsilon)x\log x-O_\varepsilon(x)$, and along a sequence with $\theta(x)\ge(A-\varepsilon)x$, $2\ge A+a-2\varepsilon-o(1)$.
So $a+A=2$.

(f) An error term $O(x)$ in a main term $2x\log x$ is a relative error $O(\frac1{\log x})$. The data say the error is $-3.67x+o(x)$, so the ratio is $1-\frac{3.67}{2\log x}$, which is
$0.886$ at $10^7$ and would reach $0.99$ only near $x=10^{80}$. There is no contradiction: the formula is precise in absolute terms, and $\log x$ grows slowly.

\tools{Tome~5, Exercise~1.6, Lemma~2.8, Proposition~1.2 and Theorem~1.3; the derivation $f\mapsto f\log$ (Tome~1, Dirichlet convolution).}

\leads Selberg's formula is the first step of the elementary proof. If $\theta(x)=x+R(x)$, it says roughly that $R(x)\log x\approx-\sum_p\log p\,R(\frac xp)$: a large error at $x$ forces
errors of the opposite sign at $\frac xp$. Selberg and Erd\H os turned this averaging into $R(x)=o(x)$, Erd\H os by showing that every interval $]x,(1+\varepsilon)x]$ has many primes.
Part (e) is the identity Selberg found in April 1948, before the full proof.
\end{solution}

\begin{solution}{pb:primroot}
\idea Being a primitive root is a condition on the index of $n$, and Möbius inversion over the divisors of $p-1$ writes its indicator as a combination of characters. The
principal character gives the expected count; every other character is controlled by Pólya--Vinogradov, uniformly over intervals.

\writeup (a) The characters with $\chi^e=\chi_0$ form a group of order $e$, the characters of $(\Z/p\Z)^*/((\Z/p\Z)^*)^e$; by orthogonality their sum at $n$ is $e$ if $n$ is an $e$-th power
and $0$ otherwise. If $n=g^k$ for a primitive root $g$, then $n$ is an $e$-th power exactly when $e\mid k$ (as $e\mid p-1$). The right side is therefore
$\sum_{e\mid\gcd(k,p-1)}\mu(e)$, which is $1$ if $\gcd(k,p-1)=1$, that is if $n$ is a primitive root, and $0$ otherwise.

(b) Sum (a) over the $n$ of an interval $I$ of $N<p$ consecutive integers (the term $n=p$, if present, is $0$ on both sides, as $\chi(p)=0$). The principal characters contribute
$\sum_{e\mid p-1}\frac{\mu(e)}e\cdot\#\{n\in I:p\nmid n\}=\frac{\varphi(p-1)}{p-1}(N-\epsilon)$ with $\epsilon\in\{0,1\}$ ($\epsilon=0$ for $I=[1,N]$). Each of the $e-1$ other characters is
primitive, since $p$ is prime, so its sum over $I$ is at most $\sqrt p(1+\log p)$ (Tome~5, Theorem~11.11). The error is at most
$\epsilon+\sum_{e\mid p-1}\mu(e)^2\frac{e-1}e\sqrt p(1+\log p)$, and $\sum\mu(e)^2\frac{e-1}e\le2^{\omega(p-1)}-\frac32$ (the term $e=1$ is $0$ and $e=2$ gives $\frac12$), which leaves room for $\epsilon\le1$.

(c) With $N=g(p)-1<p$, $P(N)=0$, so $\frac{\varphi(p-1)}{p-1}(g(p)-1)<2^{\omega(p-1)}\sqrt p(1+\log p)$. Since $2^{\omega(m)}\le\tau(m)=O(m^\varepsilon)$ (Tome~5, Exercise~2.6) and
$\frac m{\varphi(m)}=\prod_{q\mid m}\frac q{q-1}\le2^{\omega(m)}$, $g(p)=O(p^{1/2+3\varepsilon})$ for every $\varepsilon>0$.

(d) For an interval of length $H=p^{1/2+\varepsilon}<p$, the count is
\[
\frac{\varphi(p-1)}{p-1}H\Bigl(1+O\Bigl(2^{\omega(p-1)}\frac{p-1}{\varphi(p-1)}\frac{\sqrt p\log p}H\Bigr)\Bigr),
\]
and the $O$ term is
$O(p^{-\varepsilon/2})$ by (c).

(e) The bound of (c) is more than $10^4$ times $2^{\omega(p-1)}\frac{p-1}{\varphi(p-1)}$ for $p$ near $10^6$, while $g(p)\le73$. The waste is in Pólya--Vinogradov: it bounds the sum over an
interval of any length $N$ by $\sqrt p\log p$, whereas for $N$ much smaller than $\sqrt p$ the character values behave like coin tosses and the sum is about $\sqrt N$. Only
intervals longer than $\sqrt p$ can be handled this way.

\tools{Orthogonality of characters (Tome~1 and Tome~5, Chapter~3); Pólya--Vinogradov (Tome~5, Theorem~11.11); Tome~5, Exercise~2.6.}

\leads Bounds of this shape go back to Vinogradov. Burgess (1962) broke the $\sqrt p$ barrier for character sums over short intervals and obtained
$g(p)=O(p^{1/4+\varepsilon})$; under the generalised Riemann hypothesis, Shoup (1992) proved that $g(p)$ is bounded by a power of $\log p$. In the other direction, $g(p)$ is
unbounded: Problem~\ref{pb:erdosnp} shows that the least non-residue, a lower bound for $g(p)$, exceeds any fixed prime for a positive proportion of primes.
\end{solution}

\begin{solution}{pb:schur}
\idea The Fourier matrix is almost an involution: $F^2$ is $p$ times a reflection, so its eigenvalues are $\pm\sqrt p$ and $\pm i\sqrt p$ and its trace, the Gauss sum, is determined by
four multiplicities. Three equations come from $\operatorname{tr}F$, $\operatorname{tr}F^2$ and $|G|^2=p$; the fourth comes from the determinant, which Vandermonde computes exactly,
sign included. That fixes the sign of $G$, which Gauss needed four years to find.

\writeup (a) $\sum_{n=0}^{p-1}\zeta^{n^2}=1+\sum_{m\ne0}(1+\chi(m))\zeta^m=\sum_{m\ne0}\chi(m)\zeta^m$, as $\sum_{m\ne0}\zeta^m=-1$. By Tome~5, Theorem~3.5, $|G|^2=p$, and
$\overline G=\sum\chi(m)\zeta^{-m}=\chi(-1)G$, so $G^2=\chi(-1)G\overline G=\chi(-1)p$. And $\operatorname{tr}F=\sum_j\zeta^{j^2}=G$.

(b) $(F^2)_{jk}=\sum_l\zeta^{l(j+k)}=p\,[j+k\equiv0]$, so $F^2=pR$, $R^2=I$ and $F^4=p^2I$. So $F$ is diagonalisable with eigenvalues among the four fourth roots of $p^2$. Then
$p=a+b+c+d$; $\operatorname{tr}F^2=p\operatorname{tr}R=p$ (only $j=0$ is fixed by $j\mapsto-j$) gives $p(a+b)-p(c+d)=p$; and $\operatorname{tr}F=\sqrt p(a-b+ic-id)=G$. Finally
$|G|^2=p$ gives $(a-b)^2+(c-d)^2=1$.

(c) $F=(x_k^j)$ with $x_k=\zeta^k$, so $\det F=\prod_{j<k}(\zeta^k-\zeta^j)$, and $\zeta^k-\zeta^j=e(\frac{j+k}{2p})\cdot2i\sin\frac{\pi(k-j)}p$, with $\sin\frac{\pi(k-j)}p>0$. The
product of the phases is $e\bigl(\frac1{2p}\sum_{j<k}(j+k)\bigr)$, and $\sum_{j<k}(j+k)=(p-1)\sum_jj=\frac{p(p-1)^2}2$, so the phase is $e(\frac{(p-1)^2}4)=1$. Hence
$\det F=i^{p(p-1)/2}\Delta$ with $\Delta=2^{p(p-1)/2}\prod_{j<k}\sin\frac{\pi(k-j)}p>0$.

(d) The product of the eigenvalues is
\[
\sqrt p^{\,p}(-1)^bi^c(-i)^d=\sqrt p^{\,p}(-1)^bi^{c-d},
\]
so $(-1)^bi^{c-d}=i^{p(p-1)/2}$. From (b), $a+b=\frac{p+1}2$ and $c+d=\frac{p-1}2$.
If $p=4m+1$, then $a+b=2m+1$ is odd, so $a\ne b$; hence $|a-b|=1$, $c=d$, and $b\in\{m,m+1\}$. As $i^{p(p-1)/2}=i^{2m(4m+1)}=(-1)^m$, $b=m$, $a-b=1$ and $G=\sqrt p$. If $p=4m+3$,
then $c+d=2m+1$ is odd, so $|c-d|=1$, $a=b=m+1$, and $(-1)^{m+1}i^{c-d}=i^{(4m+3)(2m+1)}=i^{2m+3}=(-1)^{m+1}i$; so $c-d=1$ and $G=i\sqrt p$.

(e) Since $\chi(-1)=-1$, pairing $a$ with $p-a$ gives $T=\sum_{a<p/2}\bigl(a\chi(a)-(p-a)\chi(a)\bigr)=2U-pS$. Splitting by parity, the even $2a$ give $2\chi(2)U$ and the odd
$p-2a$ give $\sum_{a<p/2}(p-2a)\chi(p-2a)=-\chi(2)(pS-2U)$, so $T=\chi(2)(4U-pS)$. Eliminating $U$: $T=\chi(2)(2T+pS)$, so $(\chi(2)-2)T=pS$. Dirichlet's formula with
$\tau(\chi)=G=i\sqrt p$ reads $L(1,\chi)=-\frac{i\pi}{p\cdot i\sqrt p}T=-\frac{\pi T}{p^{3/2}}$. Hence $S=\frac{(2-\chi(2))\sqrt p}\pi L(1,\chi)$. Since $L(1,\chi)>0$ for a real character
(Tome~5, Theorem~3.8 and Exercise~3.4), $S>0$; as $S$ is an integer, $S\ge1$ and $L(1,\chi)\ge\frac\pi{(2-\chi(2))\sqrt p}\ge\frac\pi{3\sqrt p}$.

(f) $\chi(2)=1$ for $p\equiv7\pmod8$ ($7$, $23$, $10\,007$) and $\chi(2)=-1$ for $p\equiv3\pmod8$ ($163$, $1019$). For instance $\frac{\sqrt7}\pi1.18741=1.0000$, $3\frac{\sqrt{163}}\pi0.246069=3.0000$ and
$\frac{\sqrt{10007}}\pi2.41818=77.00$. For $p=163$, $S=3$ and $2-\chi(2)=3$, so $L(1,\chi)=\frac\pi{\sqrt{163}}$ exactly.

\tools{Tome~5, Theorems~3.5, 3.8 and~3.10, Exercise~3.4; Vandermonde determinants (Tome~2); eigenvalues and traces (Tome~2).}

\leads Gauss conjectured the sign in 1801 and proved it in 1805, after four years of trying; Schur's proof (1921) is the one given here. Dirichlet's class number formula says that
$h(-p)=\frac{\sqrt p}\pi L(1,\chi)$ for $p\equiv3\pmod4$, $p>3$, so (e) reads $h(-p)=\frac S{2-\chi(2)}$: the class number counts how many more residues than non-residues lie in the first
half. The values above are $h=1,3,1,13,77$; $h(-163)=1$ is the reason for Euler's polynomial (Problem~\ref{pb:euler41}). No proof of $S>0$ is known that avoids analysis.
\end{solution}

\begin{solution}{pb:euler41}
\idea A prime factor $p<q$ of $n^2+n+q$ is a value of a quadratic form of discriminant $1-4q$ which cannot be equivalent to $x^2+xy+qy^2$, since the latter does not represent small
numbers. So composite values and extra classes of forms are the same thing. For the density of primes, each prime $p$ divides $f(n)$ for $1+\chi(p)$ residues $n$, and $f$ is lucky:
$\chi(p)=-1$ for every $p<41$.

\writeup (a) $x\mapsto x+my$ turns $[a,b,c]$ into $[a,b+2ma,c']$, and $(x,y)\mapsto(-y,x)$ turns it into $[c,-b,a]$. Choose $m$ with $-a<b+2ma\le a$; if then $a>c$, swap. Each swap
lowers $a$, a positive integer, so the process stops at a form with $|b|\le a\le c$; if $b=-a$ translate once more, and if $a=c$ and $b<0$ swap. A reduced form has
$|D|=4ac-b^2\ge4a^2-a^2=3a^2$. A unimodular substitution maps coprime pairs bijectively to coprime pairs, which gives the last claim.

(b) Let $[a,b,c]$ be reduced of discriminant $D=1-4q$; $b$ is odd. Replacing $n$ by $-1-n$ does not change $f(n)$, so write $|b|=2n+1$ with $n\ge0$. Then
$ac=\frac{b^2-D}4=n^2+n+q=f(n)$, and $n<\frac a2\le\frac12\sqrt{(4q-1)/3}<\sqrt{q/3}\le q-2$ for $q\ge3$ (for $q=2$, $a=1$ is forced). If all $f(n)$, $0\le n\le q-2$, are prime, then
$a=1$, $|b|=1$, and $b=1$ by the convention: $[1,1,q]$ is the only reduced form. Conversely, suppose $f(n)$ is composite for some $0\le n\le q-2$. Since
$f(n)\le q^2-2q+2<q^2$, its least prime factor $p$ is $<q$. The form $[p,2n+1,\frac{f(n)}p]$ has discriminant $D$ and represents $p$ with $(x,y)=(1,0)$. But
$x^2+xy+qy^2=(x+\frac y2)^2+(q-\frac14)y^2$ is $\ge q$ when $y\ne0$, and $x^2$ with $\gcd(x,0)=1$ is $1$. So $[1,1,q]$ does not represent $p$ with coprime $(x,y)$, and by (a) the reduced form
equivalent to $[p,2n+1,\frac{f(n)}p]$ is another one.

(c) For $D=-163$: $a\le\sqrt{163/3}<8$, $b$ odd, $ac=\frac{b^2+163}4$, which is $41,43,47,53$ for $|b|=1,3,5,7$, all prime. So $a=1$, $|b|\le1$, and $[1,1,41]$ is the only reduced
form: $n^2+n+41$ is prime for $0\le n\le39$. The list $2,3,5,11,17,41$ is complete. If $1-4q=f^2D_0$ with $f>1$, then $f$ times the principal form of discriminant $D_0$ is a second
reduced form; so $D=1-4q$ must be a fundamental discriminant with one class of forms, and Heegner (1952), Baker (1966) and Stark (1967) proved that the only such $D<-4$ are
$-7,-8,-11,-19,-43,-67,-163$.

(d) For odd $p$, $4f(n)=(2n+1)^2+163$, and $2n+1$ runs over all residues modulo $p$ with $n$, so the number of roots is the number of $y$ with $y^2\equiv-163$, that is $1+\chi(p)$. And
$f(n)=n(n+1)+41$ is odd. For a random integer near $f(n)$, the chance of being prime is $\frac1{\log f(n)}$, and the chance of avoiding the prime factor $p$ is $1-\frac1p$; for $f(n)$ it is
$1-\frac{w(p)}p$ with $w(p)$ the number of roots. Correcting prime by prime gives $C=\prod_p\frac{1-w(p)/p}{1-1/p}$: the factor at $2$ is $\frac1{1/2}=2$, and at odd $p$ it is
$\frac{p-1-\chi(p)}{p-1}$. For the faster product, $\prod_{p>2}(1-\frac{\chi(p)}p)=\frac{1}{L(1,\chi)}\cdot\frac1{1-\chi(2)/2}=\frac{\sqrt{163}}\pi\cdot\frac23$, since $\chi(2)=-1$ ($-163\equiv5\pmod8$) and
$L(1,\chi)=\frac\pi{\sqrt{163}}$ (Problem~\ref{pb:schur}(f); $\chi$ is the Legendre symbol modulo $163$, by reciprocity). The Euler product of $L(1,\chi)$ converges, in the order of the primes,
to $L(1,\chi)$, by the prime number theorem in progressions. The quotients $\frac{1-\chi(p)/(p-1)}{1-\chi(p)/p}=1+O(p^{-2})$ make the remaining product absolutely convergent.

(e) For $n^2+1$, the factor at $2$ is $\frac{1-1/2}{1-1/2}=1$: half the values are even. For $f(n)=n^2+n+41$, the factor at $2$ is $2$, and $\chi(p)=-1$ for every prime $p\le37$ (otherwise
$p$ would divide some $f(n_0)$ with $0\le n_0<p$, a prime larger than $p$), so each such $p$ contributes $1+\frac1{p-1}$ instead of $1\pm\frac1{p-1}$. The primes up to $37$ alone give
a ratio $5.07$, and the larger primes, where $\chi$ and $\chi_4$ take both signs, bring it to $\frac{6.6395}{1.3728}=4.84$; the counts give $\frac{261\,081}{54\,110}=4.83$.

\tools{Binary quadratic forms (Tome~1); Tome~5, Theorem~3.10 and Theorem~5.8; Problem~\ref{pb:schur}.}

\leads Rabinowitsch proved (b) in 1913. Hardy and Littlewood stated their conjecture in 1923; it is open for every quadratic polynomial, even for $n^2+1$, which is one of
Landau's four problems of 1912. Jacobson and Williams (2003) searched for quadratic polynomials whose constant $C$ beats Euler's.
\end{solution}

\begin{solution}{pb:erdosnp}
\idea The least non-residue is $p_k$ when $p$ is a residue modulo $p_1,\dots,p_{k-1}$ and a non-residue modulo $p_k$, which by reciprocity is a condition on $p$ modulo $M_k$: half of the
classes at each step, so probability $2^{-k}$, exactly like a coin tossed until it shows tails. The difficulty is uniformity: (a) is for fixed $k$, and to average $n(p)$ one must control
large $k$, where Brun--Titchmarsh gives an upper bound good enough to sum.

\writeup (a) For $p>p_k$, $n(p)=p_k$ means $\bigl(\frac{p_i}p\bigr)=1$ for $i<k$ and $\bigl(\frac{p_k}p\bigr)=-1$. By quadratic reciprocity, $(\frac2p)$ depends on $p$ modulo $8$ and $(\frac qp)$
on $p$ modulo $4q$, so $p\mapsto\bigl((\frac{p_i}p)\bigr)_{i\le k}$ factors through a homomorphism $(\Z/M_k\Z)^*\to\{\pm1\}^k$. It is onto: for a non-empty set $I$, the product
$\prod_{i\in I}(\frac{p_i}p)$ is the Kronecker symbol of the non-square $\prod_{i\in I}p_i$, a non-principal character. So each sign pattern is taken on $\frac{\varphi(M_k)}{2^k}$ classes,
and by the prime number theorem in progressions (Tome~5, Theorem~5.8) its proportion of primes tends to $2^{-k}$.

(b) $n(p)>y$ means that all primes $q\le y$ are residues: a condition on $p$ modulo $M=8\prod_{3\le q\le y}q$, met by $\frac{\varphi(M)}{2^{\pi(y)}}$ classes. By Brun--Titchmarsh, each class
contains at most $(2+o(1))\frac x{\varphi(M)\log(x/M)}\le(4+o(1))\frac x{\varphi(M)\log x}$ primes $p\le x$, since $M\le\sqrt x$. Summing over the classes, and adding the $O(y)$ primes
dividing $M$, gives $O\bigl(\frac{\pi(x)}{2^{\pi(y)}}\bigr)$; note $y\le\log M\le\frac12\log x$ and $2^{\pi(y)}\le\sqrt x$, so $y=O(\frac{\pi(x)}{2^{\pi(y)}})$.

(c) Let $Y=\frac{\log x}3$ and fix $k_0$. The primes with $n(p)=p_k$, $k\le k_0$, contribute $\pi(x)\bigl(\sum_{k\le k_0}\frac{p_k}{2^k}+o(1)\bigr)$ by (a). If $n(p)=p_j$ with $j>k_0$, let $j'\le j$
be the least index with $p_{j'}\ge\min(p_j,Y)$; then $\min(n(p),Y)\le p_{j'}$ and $n(p)>p_{j'-1}$, with $p_{j'-1}<Y$. So these primes contribute at most
$\sum_{j>k_0,\,p_{j-1}<Y}p_j\,\#\{p\le x:n(p)>p_{j-1}\}$. For $y<Y$, $8\prod_{q\le y}q\le8e^{1.01Y}\le\sqrt x$ for large $x$ (Chebyshev), so (b) applies and the contribution is at most
$c\,\pi(x)\sum_{j>k_0}\frac{p_j}{2^{j-1}}$. This tail of a convergent series tends to $0$ as $k_0\to\infty$, which proves the limit $E$.

(d) For $x=10^8$, $Y=6.14$, so $\min(n(p),Y)$ is $2,3,5$ or $6.14$, and the limit of (c) at this $x$ is nearly $2\cdot\frac12+3\cdot\frac14+5\cdot\frac18+6.14\cdot\frac18=3.1425$: the
truncation, not a failure of the law. To reach the untruncated mean, one must show $\sum_{n(p)>\log x/3}n(p)=o(\pi(x))$. Each term is at most $p^{0.16}$ by Burgess, but this needs the
number of primes with $n(p)>y$ to be very small for $\log x\ll y\le x^{0.16}$, where (b) no longer applies. Linnik's large sieve (Tome~5, Theorem~8.10) shows that
$\#\{p\le x:n(p)>x^\varepsilon\}$ is bounded in terms of $\varepsilon$ alone, and the range in between needs a further argument.

\tools{Quadratic reciprocity (Tome~1); Tome~5, Theorem~5.8, Exercise~8.4, Theorem~8.10.}

\leads Erd\H os proved in 1961 that the untruncated mean tends to $E=\sum\frac{p_k}{2^k}$, and Martin and Pollack (2013) computed the analogous average for all non-principal characters
modulo $p$. The record value $n(p)=67$ below $10^8$ is tiny compared with $p^{0.16}$, and well below the bound $2\log^2p$ that holds under the generalised Riemann hypothesis
(Tome~5, Chapter~11).
\end{solution}

\begin{solution}{pb:hecke}
\idea Dirichlet sorted primes by residue classes with characters of $(\Z/q\Z)^*$; Hecke sorted Gaussian primes by angle with the characters $e^{4ik\theta}$ of the circle. The
proof of the prime number theorem goes through unchanged: the $L$-functions continue a little to the left of $\operatorname{Re}s=1$, do not vanish on the line by the $3$-$4$-$1$ trick,
and Newman's theorem finishes. Weyl's criterion turns ``$\vartheta_k=o(x)$ for all $k\ne0$'' into equidistribution.

\writeup (a) A unit is $i^j$ and $(i^j)^{4k}=1$, so $\lambda_k$ depends only on $(\alpha)$, and it is multiplicative because $\frac{\alpha\beta}{|\alpha\beta|}=\frac\alpha{|\alpha|}\frac\beta{|\beta|}$. Let $Q_\alpha=\alpha+[0,1[^2$.
For $|\alpha|\ge2$ and $z\in Q_\alpha$, the arguments of $z$ and $\alpha$ differ by $O(\frac1{|\alpha|})$, so $|\lambda_k(\alpha)-(\frac z{|z|})^{4k}|=O(\frac{|k|}{|\alpha|})$. Hence
\[
S_k(x)=\int_{\bigcup Q_\alpha}\Bigl(\frac z{|z|}\Bigr)^{4k}\dd A+O\Bigl(\sum_{0<|\alpha|\le\sqrt x}\frac{|k|}{|\alpha|}\Bigr)+O(1).
\]
The union of the squares differs from the disc of radius $\sqrt x$ by a set of area $O(\sqrt x)$, the integral over the disc is $\int_0^{\sqrt x}\int_0^{2\pi}e^{4ik\theta}r\dd\theta\dd r=0$,
and $\sum_{|\alpha|\le\sqrt x}\frac1{|\alpha|}=O(\sqrt x)$ by comparison with $\int\frac{\dd A}{|z|}$.

(b) $L_k(s)=\sum_na_k(n)n^{-s}$ with $\sum_{n\le x}a_k(n)=\frac14S_k(x)=O(\sqrt x)$, so the series converges and is holomorphic on $\operatorname{Re}s>\frac12$ (partial summation, Tome~5,
Lemma~1.1). Since $\Z[i]$ has unique factorisation (Tome~1, Theorem~10.3), $L_k(s)=\sum_I\lambda_k(I)N(I)^{-s}$ over non-zero ideals is an Euler product over prime ideals. These are
(Tome~1, Theorem~10.4) $(1+i)$ with $\lambda_k=(\frac{1+i}{\sqrt2})^{4k}=e^{i\pi k}=(-1)^k$; $(a+bi)$ and $(a-bi)$ for $p\equiv1\pmod4$, with $\lambda_k=e^{\pm4ik\varphi_p}$; and $(p)$, of norm
$p^2$ and $\lambda_k=1$, for $p\equiv3\pmod4$. For $k=0$, $L_0(s)=\frac14\sum r_2(n)n^{-s}=\sum(1*\chi_4)(n)n^{-s}=\zeta(s)L(s,\chi_4)$ (Tome~1, Theorem~10.5).

(c) For $\sigma>1$,
\[
\log\bigl|L_0(\sigma)^3L_k(\sigma+it)^4L_{2k}(\sigma+2it)\bigr|=\sum_{\mathfrak p,m}\frac{N\mathfrak p^{-m\sigma}}m(3+4\cos\theta+\cos2\theta)\ge0,
\]
with
$\theta=m(\arg\lambda_k(\mathfrak p)-t\log N\mathfrak p)$, because $\lambda_{2k}=\lambda_k^2$. If $L_k(1+it)=0$ with $k\ne0$, then as $\sigma\to1^+$, $L_0(\sigma)^3=O((\sigma-1)^{-3})$,
$L_k(\sigma+it)^4=O((\sigma-1)^4)$ and $L_{2k}(\sigma+2it)$ stays bounded (it is holomorphic there, as $2k\ne0$): the product tends to $0$, a contradiction.

(d) In $-\frac{L_k'}{L_k}(s)=\sum_{\mathfrak p,m}\lambda_k(\mathfrak p)^m\log N\mathfrak p\,N\mathfrak p^{-ms}$, the terms with $m=1$ and $N\mathfrak p=p\equiv1\pmod4$ give $\sum_p2\cos(4k\varphi_p)\log p\,p^{-s}$,
and the other terms form a series absolutely convergent on $\operatorname{Re}s>\frac12$. So $\Phi_k(s)=\int_1^\infty x^{-s}\dd\vartheta_k(x)$ is $-\frac{L_k'}{L_k}(s)$ plus a function holomorphic on
$\operatorname{Re}s>\frac12$. For $k\ne0$, $\Phi_k$ is holomorphic on $\operatorname{Re}s\ge1$ by (b) and (c); and $\Phi_0(s)-\frac1{s-1}$ is holomorphic there, since $\zeta(s)L(s,\chi_4)$ has a
simple pole at $1$ and no zero on the line (Tome~5, Theorems~3.8 and~5.2, and the same argument for $L(s,\chi_4)$). The functions $A_\pm=\vartheta_0\pm\vartheta_k$ are
non-decreasing (their increments are $2(1\pm\cos)\log p\ge0$) and $O(x)$, and $\int_1^\infty(A_\pm(x)-x)x^{-s-1}\dd x=\frac{\Phi_0(s)\pm\Phi_k(s)}s-\frac1{s-1}$ is holomorphic on
$\operatorname{Re}s\ge1$. By Newman's theorem $\int_1^\infty\frac{A_\pm(x)-x}{x^2}\dd x$ converges, so $A_\pm(x)\sim x$ by Lemma~5.4, and $\vartheta_k=\frac12(A_+-A_-)=o(x)$.

(e) By partial summation,
\[
\sum_{p\le x,\ p\equiv1\,(4)}\cos(k\psi_p)=o(\pi(x))\qquad(k\ne0),
\]
where $\psi_p=4\varphi_p\in\,]0,\pi[$. List the primes $p\equiv1\pmod4$ in order, and form the
sequence $\frac{\psi_{p}}{2\pi},-\frac{\psi_{p}}{2\pi},\dots$ modulo $1$, each prime giving two terms. Its Weyl sums are
\[
\frac1{2N}\sum\bigl(e^{ik\psi_p}+e^{-ik\psi_p}\bigr)=\frac1N\sum\cos(k\psi_p)\to0,
\]
so it is equidistributed. The positive copies lie in $]0,\frac12[$ and the negative ones in $]\frac12,1[$, so for $[\alpha,\beta]\subset\,]0,\pi[$ the proportion of $p$ with
$\psi_p\in[\alpha,\beta]$ tends to $2\cdot\frac{\beta-\alpha}{2\pi}=\frac{\beta-\alpha}\pi$. So $\psi_p$ is uniform on $]0,\pi[$, and $\varphi_p$ on $]0,\frac\pi4[$.

(f) $b<\frac a{10}$ means $\varphi_p<\arctan\frac1{10}$, which has predicted proportion $\frac4\pi\arctan\frac1{10}=0.12690$, against $0.12658$ observed.

\tools{Tome~1, Theorems~10.3--10.5; Tome~5, Lemma~1.1, Theorems~3.8, 5.2, 5.3 and Lemma~5.4; Weyl's criterion (Tome~4, Theorem~6.13).}

\leads Hecke introduced these characters in 1918--1920 for every imaginary quadratic field, and in general ``Grössencharaktere'' for all number fields. Fixing $b$ instead of the angle is
much harder: Friedlander and Iwaniec proved in 1998 that there are infinitely many primes $a^2+b^4$, and Heath-Brown (2001) did the same for $a^3+2b^3$. Primes $a^2+1$, that
is $b=1$, are still out of reach (Problem~\ref{pb:euler41}).
\end{solution}

\begin{solution}{pb:shifted}
\idea Each small prime $q$ divides $p-1$ with probability $\frac1{q-1}$, independently, as the prime number theorem in progressions says; so $\omega(p-1)$ behaves like a sum of independent
coins with $\sum\frac1{q-1}\approx\log\log x$, the Erd\H os--Kac picture for shifted primes. Upper bounds for large $q$ come cheaply from Brun--Titchmarsh; matching lower bounds
need Bombieri--Vinogradov.

\writeup (a) Let $Q=\prod_{q\le y}q$. For $p\nmid Q$, the set of $q\le y$ dividing $p-1$ is determined by $p$ modulo $Q$, and by the Chinese remainder theorem the classes modulo $Q$ are
tuples of classes modulo the $q$. Each class prime to $Q$ has proportion $\frac1{\varphi(Q)}$ of the primes (Tome~5, Theorem~5.8). So for $S\subset\{q\le y\}$ the proportion of primes with
$q\mid p-1$ exactly for $q\in S$ tends to $\prod_{q\in S}\frac1{q-1}\prod_{q\notin S}\frac{q-2}{q-1}$, the probability that $B_q=1$ exactly for $q\in S$.

(b) $\frac1{q-1}=\frac1q+\frac1{q(q-1)}$, so by Mertens (Tome~5, Theorem~1.6) $\sum_{q\le y}\frac1{q-1}=\log\log y+M+\sum_q\frac1{q(q-1)}+o(1)$, and
$0.261497+0.773157=1.034654$. The variance $\sum_{q\le y}\frac1{q-1}(1-\frac1{q-1})$ is $\log\log y+O(1)\to\infty$ and the $B_q$ are bounded, so Lindeberg's theorem (Tome~3,
Theorem~6.10) applies, and the difference between the mean and $\log\log y$ is $O(1)=o(\sqrt{\log\log y})$.

(c) $\sum_{p\le x}\omega(p-1)=\sum_q\pi(x;q,1)$. At most one $q>\sqrt x$ divides $p-1$, so these $q$ contribute at most $\pi(x)$. Fix $\varepsilon>0$. For $q\le x^\varepsilon$,
Brun--Titchmarsh gives $\pi(x;q,1)\le\frac{(2+o(1))x}{(q-1)(1-\varepsilon)\log x}$, uniformly; summing with (b), these $q$ give at most $\frac{2+o(1)}{1-\varepsilon}\pi(x)\log\log x$. For
$x^\varepsilon<q\le\sqrt x$, $\pi(x;q,1)\le\frac{(4+o(1))x}{(q-1)\log x}$ and $\sum_{x^\varepsilon<q\le\sqrt x}\frac1{q-1}=\log\frac1{2\varepsilon}+o(1)$, so they give $O_\varepsilon(\pi(x))$. Let $\varepsilon\to0$.

(d) With $Q=\frac{\sqrt x}{\log^Bx}$, Bombieri--Vinogradov gives $\sum_{q\le Q}\bigl|\pi(x;q,1)-\frac{\operatorname{li}(x)}{q-1}\bigr|=O(\frac x{\log^2x})$, so the $q\le Q$ contribute
$\operatorname{li}(x)\sum_{q\le Q}\frac1{q-1}+O(\frac x{\log^2x})=\pi(x)(\log\log x+O(1))$. The $q\in\,]Q,\sqrt x]$ contribute $O(\pi(x)\sum_{Q<q\le\sqrt x}\frac1q)=o(\pi(x))$ by Brun--Titchmarsh, and the
$q>\sqrt x$ at most $\pi(x)$.

(e) The model says that $q\mid p-1$ with probability $\frac1{q-1}$ for every $q<x$. For $q\le\sqrt x$ this is (d) on average, but for $q>\sqrt x$ it is a statement about primes $p$
with $\frac{p-1}q$ small and $q$ prime, a twin-prime-like problem; so the constant $1.035$ is a conjecture, which the data support: the deficits $0.228$, $0.180$, $0.143$, $0.119$,
$0.102$ shrink roughly like $\frac2{\log x}$. The variance is another matter. In the model the indicators of large $q$ are independent, so two or three large prime factors can
occur together; in reality at most one prime factor of $p-1$ exceeds $\sqrt x$, at most two exceed $x^{1/3}$, and so on. These constraints make the large-$q$ indicators negatively
correlated and remove a constant from the variance, which is of size $\log\log x+O(1)$ with a large negative $O(1)$. The same happens for $\omega(n)$: variance $1.21$ at $10^8$,
against $\log\log10^8=2.91$. A constant is a lot when $\log\log x$ is $2.9$.

\tools{Tome~5, Theorems~1.6 and~5.8, Exercise~8.4, the remark after Proposition~8.11; Lindeberg (Tome~3, Theorem~6.10).}

\leads Erd\H os proved in 1935, with Brun's sieve, that $\omega(p-1)$ has normal order $\log\log p$. The central limit theorem for $\omega(p-1)$, the analogue of Erd\H os--Kac
(Tome~5, Theorem~9.7), was proved later along the same lines, with sieve bounds for primes in progressions replacing the trivial count of multiples. The same method handles $\omega(p+a)$ for any fixed $a\ne0$.
\end{solution}

\begin{solution}{pb:romanoff}
\idea There are about $\frac x{\log x}\cdot\log_2x\asymp x$ pairs $(p,k)$ with $p+2^k\le x$, so $\mathcal R$ has positive density unless the sums collide too often. Collisions are pairs of
primes differing by $2^k-2^l$, which a sieve bounds by the twin prime count times $(\frac d{\varphi(d)})^2$. The whole difficulty is that $2^m-1$ might have unusually many small prime
factors for many $m$; it does not, because a prime $q$ divides $2^m-1$ only when $\operatorname{ord}_q(2)\mid m$, and orders are rarely small.

\writeup (a) For $p\le\frac x2$ and $2^k\le\frac x2$, $p+2^k\le x$: at least $\pi(\frac x2)\lfloor\log_2\frac x2\rfloor\ge c\,x$ pairs. By Cauchy--Schwarz,
$(\sum_{n\le x}r(n))^2\le\#\{n\le x:r(n)>0\}\cdot\sum_{n\le x}r(n)^2$.

(b) $\sum_{n\le x}r(n)^2$ counts quadruples with $p+2^k=p'+2^l\le x$. The diagonal $k=l$ gives $\sum r(n)\le\pi(x)\log_2x=O(x)$. For $k>l$, $p'=p+2^l(2^m-1)$ with $m=k-l$, and
$d=2^l(2^m-1)$ has $\frac d{\varphi(d)}=2\frac{2^m-1}{\varphi(2^m-1)}$, so the number of $p$ is at most $4c\,f(m)\frac x{\log^2x}$. Summing over $l\le\log_2x$ and $m\le\log_2x$, and
doubling for $k<l$, gives $O\bigl(\frac x{\log x}\sum_{m\le\log_2x}f(m)\bigr)=O(x)$ if $\sum_{m\le L}f(m)=O(L)$.

(c) $\frac n{\varphi(n)}=\prod_{q\mid n}\frac q{q-1}$, so $(\frac n{\varphi(n)})^2=\prod_{q\mid n}(1+g(q))=\sum_{e\mid n}\mu(e)^2g(e)$ with $1+g(q)=\frac{q^2}{(q-1)^2}$, that is
$g(q)=\frac{2q-1}{(q-1)^2}$; and $g(q)\le\frac5q$ is $3q^2-9q+5\ge0$, true for $q\ge3$. As $2^m-1$ is odd, and $e\mid2^m-1$ if and only if $\operatorname{ord}_e(2)\mid m$,
\[
\sum_{m\le L}f(m)=\sum_{e\text{ odd}}\mu(e)^2g(e)\,\#\{m\le L:\operatorname{ord}_e(2)\mid m\}\le L\sum_{e\text{ odd}}\frac{\mu(e)^2g(e)}{\operatorname{ord}_e(2)}.
\]

(d) If $\operatorname{ord}_e(2)=j\le J$, then $e\mid2^j-1\mid P_J$, and $P_J<\prod_{j\le J}2^j\le2^{J^2}$. So $A(J)=\sum_{\operatorname{ord}_e(2)\le J}\mu(e)^2g(e)\le\prod_{q\mid P_J}(1+\frac5q)$. $P_J$ has
fewer than $J^2$ prime factors, so the product is at most $\prod_{q\le p_{J^2}}(1+\frac5q)\le\exp(5\sum_{q\le p_{J^2}}\frac1q)=O((\log J)^5)$ by Mertens, as $p_{J^2}=O(J^3)$. By partial summation,
$\sum_e\frac{\mu(e)^2g(e)}{\operatorname{ord}_e(2)}=\sum_j\frac{A(j)-A(j-1)}j=\sum_j\frac{A(j)}{j(j+1)}<\infty$. With (a)--(c), $\#(\mathcal R\cap[1,x])\ge\frac{c^2x^2}{O(x)}$.

(e) The product is $2\cdot3\cdot5\cdot7\cdot13\cdot17\cdot241=11\,184\,810$, and $n\equiv7\,629\,217$ is odd. The orders of $2$ modulo $3,5,7,13,17,241$ are $2,4,3,12,8,24$, and $n\equiv1,2,1,11,8,121$ modulo these primes. So $3\mid n-2^k$ when $k\equiv0\pmod2$,
$5\mid n-2^k$ when $k\equiv1\pmod4$, $7\mid n-2^k$ when $k\equiv0\pmod3$, $13\mid n-2^k$ when $k\equiv7\pmod{12}$, $17\mid n-2^k$ when $k\equiv3\pmod8$, and $241\mid n-2^k$ when
$k\equiv23\pmod{24}$. These congruences cover all $k$: even $k$ by the first; $k\equiv1\pmod4$ by the second; $k\equiv3\pmod8$ by the fifth; and $k\equiv7\pmod8$ is $7$, $15$ or $23$ modulo
$24$, covered by the fourth, the third and the sixth. The notebook checks all $24$ residues of $k$. So $n-2^k$ is
divisible by one of the six primes, and is not prime unless it equals that prime; for each such prime this happens for at most one $k$ per $n$, and only for a sparse set of $n$
in the progression. So a positive proportion, $\frac1{11\,184\,810}$ up to sparse exceptions, of all integers are odd and not in $\mathcal R$.

(f) $n=p+2^k$ with $k\ge1$ is odd unless $p=2$, so $\mathcal R$ contains almost no even numbers, and its density is at most $\frac12$; the data are half the odd proportion. An odd $n$
has on average about $\sum_{k\le\log_2n}\frac{2}{\log n}\approx\frac2{\log2}=2.9$ representations, adjusted by local factors from $3$, $5$, $7$, \dots; with so few expected
representations, a noticeable fraction of odd numbers get none, as the $8\%$ in the data show. The covering progression of (e) is far too thin to explain them. The values of
$\frac1L\sum f(m)$ level off near $2.5$, as (c)--(d) predict.

\tools{Tome~5, Theorem~8.2 and Exercise~8.5 (Selberg's sieve), Theorem~1.3 (Mertens), Proposition~1.2; orders modulo $e$ (Tome~1).}

\leads Romanoff proved the theorem in 1934. Erd\H os found the covering of (e) in 1950, and asked whether covering systems with all moduli distinct can have arbitrarily large
least modulus; Hough proved in 2015 that they cannot. The exact density of $\mathcal R$ is unknown.
\end{solution}

\begin{solution}{pb:dedekind}
\idea A Dedekind sum is the correlation of the sawtooth with a permuted copy of itself. In the Fourier basis of $\Z/k\Z$ the sawtooth becomes cotangents, so $s(h,k)$ is a cotangent
sum; and cotangent sums are residues. The function $\cot(\pi z)\cot(\pi hz)\cot(\pi kz)$ has residues giving $s(h,k)$, $s(k,h)$ and an explicit term at $z=1$, and its integral around a
period rectangle is computable. That is the reciprocity law, which makes $s(h,k)$ computable along Euclid's algorithm.

\writeup (a) Let $w=e(\frac jk)$, $1\le j\le k-1$, and $0<r<k$. Then $\cot\frac{\pi j}k\sin\frac{2\pi jr}k=\frac{(w+1)(w^r-w^{-r})}{2(w-1)}=\frac12(w+1)\sum_{m=-r}^{r-1}w^m=\frac12\bigl(w^{-r}+w^r+2\sum_{|m|<r}w^m\bigr)$.
Summing over $j$ with $\sum_{j=1}^{k-1}w^m=k[k\mid m]-1$ gives $\frac12\bigl(-2+2(k-1)-2(2r-2)\bigr)=k-2r$, and $-\frac{k-2r}{2k}=\frac rk-\frac12=((\frac rk))$. For $r=0$ both sides vanish.
Then, using $\sum_{r\bmod k}\sin\frac{2\pi jr}k\sin\frac{2\pi j'r}k=\frac k2([j\equiv j']-[j\equiv-j'])$ and the substitution $j'\mapsto hj'$,
\[
s(h,k)=\frac1{4k^2}\sum_{j,j'}\cot\frac{\pi j}k\cot\frac{\pi j'}k\sum_r\sin\frac{2\pi jr}k\sin\frac{2\pi hj'r}k=\frac1{4k}\sum_{j=1}^{k-1}\cot\frac{\pi j}k\cot\frac{\pi hj}k,
\]
the two Kronecker terms giving equal contributions since $\cot$ is odd.

(b) $\Phi(z)=\cot(\pi z)\cot(\pi hz)\cot(\pi kz)$ has period $1$, so the vertical sides cancel. As $Y\to\infty$, $\cot(\pi(x\pm iY))\to\mp i$ uniformly, so the top side (traversed
leftwards) tends to $-(-i)^3=-i$ and the bottom side to $i^3=-i$: the integral tends to $-2i$, and the residues add up to $-\frac1\pi$. At $z=\frac jh$, $0<j<h$, only $\cot(\pi hz)$ has
a pole, with residue $\frac1{\pi h}$; the residues there add up to $\frac1{\pi h}\sum_j\cot\frac{\pi j}h\cot\frac{\pi kj}h=\frac4\pi s(k,h)$, and those at $\frac jk$ to $\frac4\pi s(h,k)$. At $z=1$,
with $\cot(\pi u)=\frac1{\pi u}-\frac{\pi u}3+O(u^3)$, $\Phi(1+u)=\frac1{\pi^3hku^3}(1-\frac{\pi^2u^2}3)(1-\frac{\pi^2h^2u^2}3)(1-\frac{\pi^2k^2u^2}3)+\dots$, with residue
$-\frac{1+h^2+k^2}{3\pi hk}$. So $\frac4\pi(s(h,k)+s(k,h))-\frac{1+h^2+k^2}{3\pi hk}=-\frac1\pi$, which is the law.

(c) $s(1,k)=\sum_{r=1}^{k-1}(\frac rk-\frac12)^2=\frac{(k-1)(2k-1)}{6k}-\frac{k-1}2+\frac{k-1}4=\frac{(k-1)(k-2)}{12k}$. Writing $\{\frac{hr}k\}=\frac{hr\bmod k}k$ and using
$\sum_{r=1}^{k-1}\{\frac{hr}k\}=\frac{k-1}2$, one finds $s(h,k)=\frac1{k^2}\sum_rr\,(hr\bmod k)-\frac{k-1}4$. Since $r(hr\bmod k)\equiv hr^2\pmod k$, $6\sum_rr(hr\bmod k)\equiv h(k-1)k(2k-1)\equiv0\pmod k$,
so $6k\,s(h,k)=\frac6k\sum_rr(hr\bmod k)-\frac{3k(k-1)}2$ is an integer. Finally, by Cauchy--Schwarz and because $r\mapsto hr$ permutes the residues,
$|s(h,k)|\le\sum_r((\frac rk))^2=s(1,k)$.

(d) $s(h,k)$ depends on $h$ modulo $k$. Let $0<h<k$ and $k=a_1h+r_1$, so that $\frac{r_1}h=[0;a_2,\dots,a_n]$. By (b),
$|s(h,k)|\le|s(r_1,h)|+\frac1{12}(\frac kh+\frac hk+\frac1{hk})+\frac14\le|s(r_1,h)|+\frac{a_1}{12}+1$, as $\frac kh<a_1+1$, $\frac hk<1$ and $\frac1{hk}\le\frac12$. Iterating down to $s(0,1)=0$ gives
$|s(h,k)|\le\frac1{12}\sum a_i+n$. The same recursion, with signs, computes $s(h,k)$ exactly, and Euclid's algorithm has $n=O(\log k)$ steps.

(e) If $\frac{s(h,k)}{\log k}$ were exactly Cauchy, $P(|s|>t\log k)\approx\frac c t$, and the mean of $|s|/\log k$ would diverge logarithmically. Since $|s(h,k)|\le s(1,k)<\frac k{12}$, the tail is cut
at $t\approx\frac k{\log k}$, and $\mathbb E\frac{|s|}{\log k}\approx\int_1^{k/\log k}\frac c t\dd t\approx c\log k$. So the average of $|s|$ grows like $\log^2k$, as the data show: the ratios to
$\log^2k$ are $0.045$, $0.043$, $0.044$, $0.044$. The large values come from $h$ with one large partial quotient, as (d) suggests: $h=1$, $h=2$, or $h$ close to $\frac ka$.

\tools{Discrete Fourier analysis on $\Z/k\Z$ (Tome~4); residues (Tome~4, Chapter~9); continued fractions and Euclid's algorithm (Tome~1).}

\leads Dedekind found the sums in 1892 in the transformation law of $\log\eta(\tau)$ under $\tau\mapsto\frac{a\tau+b}{c\tau+d}$, which involves $s(d,c)$; Tome~5, Theorem~12.7 is the case
$s(0,1)=0$. The proof of reciprocity given here is one of several in Rademacher and Grosswald's monograph.
\end{solution}

\section*{Map of the chapter}

\begin{center}
\small
\adjustbox{max width=\linewidth}{\begin{tabular}{lll}
Problem & Topic & Tomes used\\\hline
\ref{pb:jacobi} & Jacobi's four-square theorem, $\zeta(2)$ by lattice points & 0, 1, 4\\
\ref{pb:densities} & local densities and Siegel's formula for $r_4(n)$ & 1\\
\ref{pb:twosquares} & how many integers are sums of two squares & 1, 4, 5\\
\ref{pb:ramanujanE2} & $E_2$ and Ramanujan's differential equations & 5\\
\ref{pb:landauzeros} & Landau's formula: the zeros hear the primes & 4, 5\\
\ref{pb:hearzeros} & the Guinand--Weil formula: the primes hear the zeros & 4, 5\\
\ref{pb:cramersign} & Cramér's coin and the sign of $\pi(x)-\operatorname{li}(x)$ & 3, 5\\
\ref{pb:selberg} & Selberg's formula and $a+A=2$ & 1, 5\\
\ref{pb:primroot} & the least primitive root & 1, 5\\
\ref{pb:schur} & Schur's determinant, the sign of the Gauss sum & 2, 5\\
\ref{pb:euler41} & Euler's polynomial, Rabinowitsch, Hardy--Littlewood & 1, 5\\
\ref{pb:erdosnp} & the least non-residue on average & 1, 5\\
\ref{pb:hecke} & Hecke's angles of the Gaussian primes & 1, 4, 5\\
\ref{pb:shifted} & the prime factors of $p-1$ & 3, 5\\
\ref{pb:romanoff} & Romanoff's theorem and Erd\H os's covering & 1, 5\\
\ref{pb:dedekind} & Dedekind sums and reciprocity & 1, 4, 5\\
\end{tabular}}
\end{center}

\section*{Sources and further reading}

\begin{itemize}
\item G.~H.~Hardy, \emph{A Mathematician's Apology}, Cambridge University Press, 1940, \S10 (the epigraph).
\item D.~Goldfeld, The elementary proof of the prime number theorem: an historical perspective, in D.~Chudnovsky, G.~Chudnovsky and M.~Nathanson (eds.), \emph{Number Theory: New York
Seminar 2003}, Springer, 2004 (the story of 1948, with the letters quoted); A.~Selberg, An elementary proof of the prime-number theorem, \emph{Annals of Mathematics} 50 (1949),
305--313; P.~Erd\H os, On a new method in elementary number theory which leads to an elementary proof of the prime number theorem, \emph{Proceedings of the National Academy of
Sciences} 35 (1949), 374--384.
\item M.~D.~Hirschhorn, A simple proof of Jacobi's four-square theorem, \emph{Proceedings of the AMS} 101 (1987), 436--438; C.~G.~J.~Jacobi, \emph{Fundamenta nova theoriae functionum
ellipticarum}, Königsberg, 1829.
\item C.~L.~Siegel, Über die analytische Theorie der quadratischen Formen, \emph{Annals of Mathematics} 36 (1935), 527--606.
\item E.~Landau, Über die Einteilung der positiven ganzen Zahlen in vier Klassen nach der Mindestzahl der zu ihrer additiven Zusammensetzung erforderlichen Quadrate, \emph{Archiv der
Mathematik und Physik} 13 (1908), 305--312.
\item S.~Ramanujan, On certain arithmetical functions, \emph{Transactions of the Cambridge Philosophical Society} 22 (1916), 159--184.
\item E.~Landau, Über die Nullstellen der Zetafunktion, \emph{Mathematische Annalen} 71 (1911), 548--564; S.~M.~Gonek, An explicit formula of Landau and its applications to the theory of the
zeta-function, \emph{Contemporary Mathematics} 143 (1993), 395--413; H.~Rademacher, Fourier analysis in number theory, Symposium on Harmonic Analysis and Related Integral Transforms,
Cornell University, 1956, in \emph{Collected Papers}, Vol.~II, MIT Press, 1974, 434--458; E.~Hlawka, Über die Gleichverteilung gewisser Folgen, welche mit den Nullstellen der
Zetafunktion zusammenhängen, \emph{Sitzungsberichte der Österreichischen Akademie der Wissenschaften, Abt.~II} 184 (1975), 459--471.
\item A.~P.~Guinand, A summation formula in the theory of prime numbers, \emph{Proceedings of the London Mathematical Society} 50 (1948), 107--119; A.~Weil, Sur les ``formules explicites'' de la
théorie des nombres premiers, \emph{Communications du Séminaire Mathématique de l'Université de Lund}, tome supplémentaire (1952), 252--265.
\item J.~E.~Littlewood, Sur la distribution des nombres premiers, \emph{Comptes Rendus de l'Académie des Sciences} 158 (1914), 1869--1872; J.~Büthe, An analytic method for bounding
$\psi(x)$, \emph{Mathematics of Computation} 87 (2018), 1991--2009; M.~Rubinstein and P.~Sarnak, Chebyshev's bias, \emph{Experimental Mathematics} 3 (1994), 173--197.
\item D.~A.~Burgess, On character sums and primitive roots, \emph{Proceedings of the London Mathematical Society} 12 (1962), 179--192; V.~Shoup, Searching for primitive roots in finite
fields, \emph{Mathematics of Computation} 58 (1992), 369--380; Encyclopaedia of Mathematics, article ``Distribution of power residues and non-residues''.
\item I.~Schur, Über die Gaussschen Summen, \emph{Nachrichten der Gesellschaft der Wissenschaften zu Göttingen} (1921), 147--153; B.~C.~Berndt, R.~J.~Evans and K.~S.~Williams, \emph{Gauss
and Jacobi Sums}, Wiley, 1998 (the history of the sign).
\item G.~Rabinowitsch, Eindeutigkeit der Zerlegung in Primzahlfaktoren in quadratischen Zahlkörpern, \emph{Journal für die reine und angewandte Mathematik} 142 (1913), 153--164;
G.~H.~Hardy and J.~E.~Littlewood, Some problems of ``Partitio numerorum'' III, \emph{Acta Mathematica} 44 (1923), 1--70; M.~J.~Jacobson and H.~C.~Williams, New quadratic polynomials
with high densities of prime values, \emph{Mathematics of Computation} 72 (2003), 499--519.
\item P.~Erd\H os, Remarks on number theory I, \emph{Matematikai Lapok} 12 (1961), 10--17; G.~Martin and P.~Pollack, The average least character non-residue and further variations on a
theme of Erd\H os, \emph{Journal of the London Mathematical Society} 87 (2013), 22--42.
\item E.~Hecke, Eine neue Art von Zetafunktionen und ihre Beziehungen zur Verteilung der Primzahlen II, \emph{Mathematische Zeitschrift} 6 (1920), 11--51; J.~Friedlander and H.~Iwaniec,
The polynomial $X^2+Y^4$ captures its primes, \emph{Annals of Mathematics} 148 (1998), 945--1040; D.~R.~Heath-Brown, Primes represented by $x^3+2y^3$, \emph{Acta Mathematica} 186 (2001),
1--84.
\item P.~Erd\H os, On the normal number of prime factors of $p-1$ and some related problems concerning Euler's $\varphi$-function, \emph{Quarterly Journal of Mathematics (Oxford)} 6
(1935), 205--213.
\item N.~P.~Romanoff, Über einige Sätze der additiven Zahlentheorie, \emph{Mathematische Annalen} 109 (1934), 668--678; P.~Erd\H os, On integers of the form $2^k+p$ and some related
problems, \emph{Summa Brasiliensis Mathematicae} 2 (1950), 113--123; R.~Hough, Solution of the minimum modulus problem for covering systems, \emph{Annals of Mathematics} 181 (2015),
361--382.
\item R.~Dedekind, Erläuterungen zu den Fragmenten XXVIII, in B.~Riemann, \emph{Gesammelte mathematische Werke}, 2nd edition, 1892; H.~Rademacher and E.~Grosswald, \emph{Dedekind Sums},
Carus Mathematical Monographs 16, Mathematical Association of America, 1972; I.~Vardi, Dedekind sums have a limiting distribution, \emph{International Mathematics Research Notices}
1993, 1--12.
\end{itemize}
